% Options for packages loaded elsewhere
% Options for packages loaded elsewhere
\PassOptionsToPackage{unicode}{hyperref}
\PassOptionsToPackage{hyphens}{url}
%
\documentclass[
  english,
  letterpaper,
]{book}
\usepackage{xcolor}
\usepackage{amsmath,amssymb}
\setcounter{secnumdepth}{5}
\usepackage{iftex}
\ifPDFTeX
  \usepackage[T1]{fontenc}
  \usepackage[utf8]{inputenc}
  \usepackage{textcomp} % provide euro and other symbols
\else % if luatex or xetex
  \usepackage{unicode-math} % this also loads fontspec
  \defaultfontfeatures{Scale=MatchLowercase}
  \defaultfontfeatures[\rmfamily]{Ligatures=TeX,Scale=1}
\fi
\usepackage{lmodern}
\ifPDFTeX\else
  % xetex/luatex font selection
\fi
% Use upquote if available, for straight quotes in verbatim environments
\IfFileExists{upquote.sty}{\usepackage{upquote}}{}
\IfFileExists{microtype.sty}{% use microtype if available
  \usepackage[]{microtype}
  \UseMicrotypeSet[protrusion]{basicmath} % disable protrusion for tt fonts
}{}
\makeatletter
\@ifundefined{KOMAClassName}{% if non-KOMA class
  \IfFileExists{parskip.sty}{%
    \usepackage{parskip}
  }{% else
    \setlength{\parindent}{0pt}
    \setlength{\parskip}{6pt plus 2pt minus 1pt}}
}{% if KOMA class
  \KOMAoptions{parskip=half}}
\makeatother
% Make \paragraph and \subparagraph free-standing
\makeatletter
\ifx\paragraph\undefined\else
  \let\oldparagraph\paragraph
  \renewcommand{\paragraph}{
    \@ifstar
      \xxxParagraphStar
      \xxxParagraphNoStar
  }
  \newcommand{\xxxParagraphStar}[1]{\oldparagraph*{#1}\mbox{}}
  \newcommand{\xxxParagraphNoStar}[1]{\oldparagraph{#1}\mbox{}}
\fi
\ifx\subparagraph\undefined\else
  \let\oldsubparagraph\subparagraph
  \renewcommand{\subparagraph}{
    \@ifstar
      \xxxSubParagraphStar
      \xxxSubParagraphNoStar
  }
  \newcommand{\xxxSubParagraphStar}[1]{\oldsubparagraph*{#1}\mbox{}}
  \newcommand{\xxxSubParagraphNoStar}[1]{\oldsubparagraph{#1}\mbox{}}
\fi
\makeatother

\usepackage{color}
\usepackage{fancyvrb}
\newcommand{\VerbBar}{|}
\newcommand{\VERB}{\Verb[commandchars=\\\{\}]}
\DefineVerbatimEnvironment{Highlighting}{Verbatim}{commandchars=\\\{\}}
% Add ',fontsize=\small' for more characters per line
\usepackage{framed}
\definecolor{shadecolor}{RGB}{241,243,245}
\newenvironment{Shaded}{\begin{snugshade}}{\end{snugshade}}
\newcommand{\AlertTok}[1]{\textcolor[rgb]{0.68,0.00,0.00}{#1}}
\newcommand{\AnnotationTok}[1]{\textcolor[rgb]{0.37,0.37,0.37}{#1}}
\newcommand{\AttributeTok}[1]{\textcolor[rgb]{0.40,0.45,0.13}{#1}}
\newcommand{\BaseNTok}[1]{\textcolor[rgb]{0.68,0.00,0.00}{#1}}
\newcommand{\BuiltInTok}[1]{\textcolor[rgb]{0.00,0.23,0.31}{#1}}
\newcommand{\CharTok}[1]{\textcolor[rgb]{0.13,0.47,0.30}{#1}}
\newcommand{\CommentTok}[1]{\textcolor[rgb]{0.37,0.37,0.37}{#1}}
\newcommand{\CommentVarTok}[1]{\textcolor[rgb]{0.37,0.37,0.37}{\textit{#1}}}
\newcommand{\ConstantTok}[1]{\textcolor[rgb]{0.56,0.35,0.01}{#1}}
\newcommand{\ControlFlowTok}[1]{\textcolor[rgb]{0.00,0.23,0.31}{\textbf{#1}}}
\newcommand{\DataTypeTok}[1]{\textcolor[rgb]{0.68,0.00,0.00}{#1}}
\newcommand{\DecValTok}[1]{\textcolor[rgb]{0.68,0.00,0.00}{#1}}
\newcommand{\DocumentationTok}[1]{\textcolor[rgb]{0.37,0.37,0.37}{\textit{#1}}}
\newcommand{\ErrorTok}[1]{\textcolor[rgb]{0.68,0.00,0.00}{#1}}
\newcommand{\ExtensionTok}[1]{\textcolor[rgb]{0.00,0.23,0.31}{#1}}
\newcommand{\FloatTok}[1]{\textcolor[rgb]{0.68,0.00,0.00}{#1}}
\newcommand{\FunctionTok}[1]{\textcolor[rgb]{0.28,0.35,0.67}{#1}}
\newcommand{\ImportTok}[1]{\textcolor[rgb]{0.00,0.46,0.62}{#1}}
\newcommand{\InformationTok}[1]{\textcolor[rgb]{0.37,0.37,0.37}{#1}}
\newcommand{\KeywordTok}[1]{\textcolor[rgb]{0.00,0.23,0.31}{\textbf{#1}}}
\newcommand{\NormalTok}[1]{\textcolor[rgb]{0.00,0.23,0.31}{#1}}
\newcommand{\OperatorTok}[1]{\textcolor[rgb]{0.37,0.37,0.37}{#1}}
\newcommand{\OtherTok}[1]{\textcolor[rgb]{0.00,0.23,0.31}{#1}}
\newcommand{\PreprocessorTok}[1]{\textcolor[rgb]{0.68,0.00,0.00}{#1}}
\newcommand{\RegionMarkerTok}[1]{\textcolor[rgb]{0.00,0.23,0.31}{#1}}
\newcommand{\SpecialCharTok}[1]{\textcolor[rgb]{0.37,0.37,0.37}{#1}}
\newcommand{\SpecialStringTok}[1]{\textcolor[rgb]{0.13,0.47,0.30}{#1}}
\newcommand{\StringTok}[1]{\textcolor[rgb]{0.13,0.47,0.30}{#1}}
\newcommand{\VariableTok}[1]{\textcolor[rgb]{0.07,0.07,0.07}{#1}}
\newcommand{\VerbatimStringTok}[1]{\textcolor[rgb]{0.13,0.47,0.30}{#1}}
\newcommand{\WarningTok}[1]{\textcolor[rgb]{0.37,0.37,0.37}{\textit{#1}}}

\usepackage{longtable,booktabs,array}
\usepackage{calc} % for calculating minipage widths
% Correct order of tables after \paragraph or \subparagraph
\usepackage{etoolbox}
\makeatletter
\patchcmd\longtable{\par}{\if@noskipsec\mbox{}\fi\par}{}{}
\makeatother
% Allow footnotes in longtable head/foot
\IfFileExists{footnotehyper.sty}{\usepackage{footnotehyper}}{\usepackage{footnote}}
\makesavenoteenv{longtable}
\usepackage{graphicx}
\makeatletter
\newsavebox\pandoc@box
\newcommand*\pandocbounded[1]{% scales image to fit in text height/width
  \sbox\pandoc@box{#1}%
  \Gscale@div\@tempa{\textheight}{\dimexpr\ht\pandoc@box+\dp\pandoc@box\relax}%
  \Gscale@div\@tempb{\linewidth}{\wd\pandoc@box}%
  \ifdim\@tempb\p@<\@tempa\p@\let\@tempa\@tempb\fi% select the smaller of both
  \ifdim\@tempa\p@<\p@\scalebox{\@tempa}{\usebox\pandoc@box}%
  \else\usebox{\pandoc@box}%
  \fi%
}
% Set default figure placement to htbp
\def\fps@figure{htbp}
\makeatother


% definitions for citeproc citations
\NewDocumentCommand\citeproctext{}{}
\NewDocumentCommand\citeproc{mm}{%
  \begingroup\def\citeproctext{#2}\cite{#1}\endgroup}
\makeatletter
 % allow citations to break across lines
 \let\@cite@ofmt\@firstofone
 % avoid brackets around text for \cite:
 \def\@biblabel#1{}
 \def\@cite#1#2{{#1\if@tempswa , #2\fi}}
\makeatother
\newlength{\cslhangindent}
\setlength{\cslhangindent}{1.5em}
\newlength{\csllabelwidth}
\setlength{\csllabelwidth}{3em}
\newenvironment{CSLReferences}[2] % #1 hanging-indent, #2 entry-spacing
 {\begin{list}{}{%
  \setlength{\itemindent}{0pt}
  \setlength{\leftmargin}{0pt}
  \setlength{\parsep}{0pt}
  % turn on hanging indent if param 1 is 1
  \ifodd #1
   \setlength{\leftmargin}{\cslhangindent}
   \setlength{\itemindent}{-1\cslhangindent}
  \fi
  % set entry spacing
  \setlength{\itemsep}{#2\baselineskip}}}
 {\end{list}}
\usepackage{calc}
\newcommand{\CSLBlock}[1]{\hfill\break\parbox[t]{\linewidth}{\strut\ignorespaces#1\strut}}
\newcommand{\CSLLeftMargin}[1]{\parbox[t]{\csllabelwidth}{\strut#1\strut}}
\newcommand{\CSLRightInline}[1]{\parbox[t]{\linewidth - \csllabelwidth}{\strut#1\strut}}
\newcommand{\CSLIndent}[1]{\hspace{\cslhangindent}#1}

\ifLuaTeX
\usepackage[bidi=basic]{babel}
\else
\usepackage[bidi=default]{babel}
\fi
% get rid of language-specific shorthands (see #6817):
\let\LanguageShortHands\languageshorthands
\def\languageshorthands#1{}
\ifLuaTeX
  \usepackage[english]{selnolig} % disable illegal ligatures
\fi


\setlength{\emergencystretch}{3em} % prevent overfull lines

\providecommand{\tightlist}{%
  \setlength{\itemsep}{0pt}\setlength{\parskip}{0pt}}



 


\usepackage{booktabs}

\usepackage{tcolorbox}
\usepackage[T1]{fontenc}

\definecolor{dangcol}{RGB}{152,62,130}
\definecolor{warncol}{RGB}{245,220,112}
\definecolor{infocol}{RGB}{70,122,172}
\definecolor{trycol}{RGB}{97,88,156}

\newtcolorbox{dangerous}{
  colback=dangcol!10,
  colframe=dangcol,
  coltext=black,
  boxsep=5pt,
  arc=4pt
}

\newtcolorbox{warning}{
  colback=warncol!10,
  colframe=warncol,
  coltext=black,
  boxsep=5pt,
  arc=4pt
}

\newtcolorbox{info}{
  colback=infocol!10,
  colframe=infocol,
  coltext=black,
  boxsep=5pt,
  arc=4pt
}

\newtcolorbox{try}{
  colback=trycol!10,
  colframe=trycol,
  coltext=black,
  boxsep=5pt,
  arc=4pt
}
\makeatletter
\@ifpackageloaded{tcolorbox}{}{\usepackage[skins,breakable]{tcolorbox}}
\@ifpackageloaded{fontawesome5}{}{\usepackage{fontawesome5}}
\definecolor{quarto-callout-color}{HTML}{909090}
\definecolor{quarto-callout-note-color}{HTML}{0758E5}
\definecolor{quarto-callout-important-color}{HTML}{CC1914}
\definecolor{quarto-callout-warning-color}{HTML}{EB9113}
\definecolor{quarto-callout-tip-color}{HTML}{00A047}
\definecolor{quarto-callout-caution-color}{HTML}{FC5300}
\definecolor{quarto-callout-color-frame}{HTML}{acacac}
\definecolor{quarto-callout-note-color-frame}{HTML}{4582ec}
\definecolor{quarto-callout-important-color-frame}{HTML}{d9534f}
\definecolor{quarto-callout-warning-color-frame}{HTML}{f0ad4e}
\definecolor{quarto-callout-tip-color-frame}{HTML}{02b875}
\definecolor{quarto-callout-caution-color-frame}{HTML}{fd7e14}
\makeatother
\makeatletter
\@ifpackageloaded{bookmark}{}{\usepackage{bookmark}}
\makeatother
\makeatletter
\@ifpackageloaded{caption}{}{\usepackage{caption}}
\AtBeginDocument{%
\ifdefined\contentsname
  \renewcommand*\contentsname{Table of contents}
\else
  \newcommand\contentsname{Table of contents}
\fi
\ifdefined\listfigurename
  \renewcommand*\listfigurename{List of Figures}
\else
  \newcommand\listfigurename{List of Figures}
\fi
\ifdefined\listtablename
  \renewcommand*\listtablename{List of Tables}
\else
  \newcommand\listtablename{List of Tables}
\fi
\ifdefined\figurename
  \renewcommand*\figurename{Figure}
\else
  \newcommand\figurename{Figure}
\fi
\ifdefined\tablename
  \renewcommand*\tablename{Table}
\else
  \newcommand\tablename{Table}
\fi
}
\@ifpackageloaded{float}{}{\usepackage{float}}
\floatstyle{ruled}
\@ifundefined{c@chapter}{\newfloat{codelisting}{h}{lop}}{\newfloat{codelisting}{h}{lop}[chapter]}
\floatname{codelisting}{Listing}
\newcommand*\listoflistings{\listof{codelisting}{List of Listings}}
\makeatother
\makeatletter
\makeatother
\makeatletter
\@ifpackageloaded{caption}{}{\usepackage{caption}}
\@ifpackageloaded{subcaption}{}{\usepackage{subcaption}}
\makeatother
\usepackage{bookmark}
\IfFileExists{xurl.sty}{\usepackage{xurl}}{} % add URL line breaks if available
\urlstyle{same}
\hypersetup{
  pdftitle={Intro to Data Visualisation Using R and the Slave Voyages dataset},
  pdfauthor={Emily Nordmann},
  pdflang={en},
  hidelinks,
  pdfcreator={LaTeX via pandoc}}


\title{Intro to Data Visualisation Using R and the Slave Voyages
dataset}
\author{Emily Nordmann}
\date{}
\begin{document}
\frontmatter
\maketitle

\renewcommand*\contentsname{Table of contents}
{
\setcounter{tocdepth}{2}
\tableofcontents
}

\mainmatter
\bookmarksetup{startatroot}

\chapter*{Overview}\label{overview}
\addcontentsline{toc}{chapter}{Overview}

\markboth{Overview}{Overview}

This is a beginner's guide to data visualisation in R using
\texttt{ggplot2}. It assumes you have installed R and RStudio and know
the absolute basics, such as how to run a line of code and load a
package. If you are starting from scratch, see
Appendix~\ref{sec-installing-r}.

The dataset used throughout is a subset of the
\href{https://www.slavevoyages.org/}{Trans-Atlantic Slave Trade
Database}, curated by the Slave Voyages project. The full database
records 36,000 individual slaving voyages between 1514 and 1866. The
subset used here focuses on voyages that originated in Britain. The
history these data describe is brutal. The figures you will produce
show, among other things, mortality rates, age and sex composition of
the people forcibly transported, and how the trade scaled with British
involvement over time. The tutorial treats the data with the seriousness
the subject demands.

\section*{Intended learning outcomes}\label{intended-learning-outcomes}
\addcontentsline{toc}{section}{Intended learning outcomes}

\markright{Intended learning outcomes}

By the end of this book you will be able to:

\begin{itemize}
\tightlist
\item
  Load tabular data into R with \texttt{read\_csv()} and inspect it
\item
  Clean variable names and types using \texttt{clean\_names()},
  \texttt{rename()}, \texttt{factor()} and \texttt{mutate()}
\item
  Build plots in layers using the \texttt{ggplot2} grammar of graphics
\item
  Choose an appropriate plot type for the data at hand
\item
  Customise axes, scales, themes and colours
\item
  Reshape data with \texttt{pivot\_longer()} to plot grouped
  distributions
\item
  Combine plots with \texttt{patchwork}
\item
  Save plots in publication-ready formats with alt text
\end{itemize}

\section*{Citing}\label{citing}
\addcontentsline{toc}{section}{Citing}

\markright{Citing}

If you use this tutorial in teaching or in your own writing, please cite
it as:

Nordmann, E. (2026). \emph{Intro to Data Visualisation Using R and the
Slave Voyages dataset.}
\url{https://github.com/emilynordmann/slavevoyagedata}

If you want a fuller treatment of data visualisation aimed at
researchers, see Nordmann et al. (2022), the published article that
accompanies the introdataviz tutorial.

\section*{Acknowledgement}\label{acknowledgement}
\addcontentsline{toc}{section}{Acknowledgement}

\markright{Acknowledgement}

This tutorial was originally written in 2020 at the request of
\href{https://twitter.com/DavidTree}{David Tree} for anti-racism
curriculum materials. It has been revised for Quarto and updated to
current tidyverse practice. The code is available under a Creative
Commons Attribution-ShareAlike licence on
\href{https://github.com/emilynordmann/slavevoyagedata}{GitHub} and you
are welcome to adapt it for your own teaching without further
permission.

\bookmarksetup{startatroot}

\chapter{Getting started}\label{sec-getting-started}

\section{Intended learning outcomes}\label{intended-learning-outcomes-1}

By the end of this chapter you will be able to:

\begin{itemize}
\tightlist
\item
  Load the \texttt{tidyverse}, \texttt{janitor} and \texttt{patchwork}
  packages
\item
  Read a CSV file into R with \texttt{read\_csv()} and verify the import
  worked
\item
  Standardise messy column names with \texttt{clean\_names()} from the
  \texttt{janitor} package
\item
  Rename variables with \texttt{rename()} to make them easier to work
  with
\item
  Convert columns to factors with \texttt{factor()}, controlling the
  display order
\item
  Use \texttt{mutate()} and \texttt{across()} to transform several
  columns at once
\end{itemize}

\section{Functions used}\label{functions-used}

\begin{itemize}
\tightlist
\item
  tidyverse (load with \texttt{library(tidyverse)})

  \begin{itemize}
  \tightlist
  \item
    readr: \texttt{read\_csv()}
  \item
    dplyr: \texttt{rename()}, \texttt{mutate()}, \texttt{across()}
  \item
    tibble: \texttt{glimpse()}
  \item
    purrr: \texttt{map\_chr()}
  \end{itemize}
\item
  janitor: \texttt{clean\_names()}
\item
  base R: \texttt{summary()}, \texttt{str()}, \texttt{factor()},
  \texttt{c()}
\end{itemize}

\section{Set-up}\label{set-up}

For this tutorial you will find it helpful to write your code in a
\href{https://quarto.org/}{Quarto} document. Quarto lets you intersperse
code, output and notes in a single file that is easy to share. If you
have not used Quarto before, the
\href{https://psyteachr.github.io/ads-v3/02-reports.html}{Quarto
Reports} chapter of \texttt{ads-v3} is a good starting point.

\begin{enumerate}
\def\labelenumi{\arabic{enumi}.}
\tightlist
\item
  Open RStudio.
\item
  Create a new RStudio project in a fresh folder.
\item
  Download this data into that folder.
\item
  Create a new Quarto document.
\item
  Replace the default setup chunk at the top of your \texttt{.qmd} file
  with the one below.
\end{enumerate}

\begin{Shaded}
\begin{Highlighting}[]
\InformationTok{\textasciigrave{}\textasciigrave{}\textasciigrave{}\{r\}}
\CommentTok{\#| label: setup}
\CommentTok{\#| include: false}

\FunctionTok{library}\NormalTok{(tidyverse)}
\FunctionTok{library}\NormalTok{(janitor)}
\FunctionTok{library}\NormalTok{(patchwork)}
\InformationTok{\textasciigrave{}\textasciigrave{}\textasciigrave{}}
\end{Highlighting}
\end{Shaded}

\section{Loading packages}\label{loading-packages}

The main package we will use is the \texttt{tidyverse} collection. This
loads \texttt{readr} for importing data, \texttt{dplyr} for wrangling,
\texttt{tidyr} for reshaping, and \texttt{ggplot2} for plotting, along
with several others. We will also use \texttt{janitor} to tidy up the
column names, and \texttt{patchwork} later to combine multiple plots
into a single figure. \texttt{janitor} and \texttt{patchwork} are
separate packages, so install them first with
\texttt{install.packages(c("janitor",\ "patchwork"))} if you have not
already.

\begin{Shaded}
\begin{Highlighting}[]
\FunctionTok{library}\NormalTok{(tidyverse)}
\FunctionTok{library}\NormalTok{(janitor)}
\FunctionTok{library}\NormalTok{(patchwork)}
\end{Highlighting}
\end{Shaded}

\begin{tcolorbox}[enhanced jigsaw, colframe=quarto-callout-tip-color-frame, breakable, left=2mm, leftrule=.75mm, toptitle=1mm, rightrule=.15mm, opacitybacktitle=0.6, arc=.35mm, colback=white, title=\textcolor{quarto-callout-tip-color}{\faLightbulb}\hspace{0.5em}{Tip}, coltitle=black, colbacktitle=quarto-callout-tip-color!10!white, toprule=.15mm, bottomtitle=1mm, opacityback=0, titlerule=0mm, bottomrule=.15mm]

You only need to install a package once
(\texttt{install.packages("tidyverse")}) but you need to load it with
\texttt{library()} every time you start a new R session.

\end{tcolorbox}

\section{Loading the data}\label{loading-the-data}

The data are in a CSV file called \texttt{data.csv}. We read them in
with \texttt{read\_csv()} from \texttt{readr}. The result is stored in
an object called \texttt{dat}.

\begin{Shaded}
\begin{Highlighting}[]
\NormalTok{dat }\OtherTok{\textless{}{-}} \FunctionTok{read\_csv}\NormalTok{(}\StringTok{"data.csv"}\NormalTok{)}
\end{Highlighting}
\end{Shaded}

When \texttt{read\_csv()} runs successfully it prints a column
specification telling you what type of data it found in each column.
\texttt{chr} means character (text), \texttt{dbl} means double (decimal
number). If you get the error message
\texttt{…does\ not\ exist\ in\ current\ working\ directory} you have
almost certainly mis-typed the file name, forgotten the \texttt{.csv}
extension, or saved the data outside your project folder. If you get
\texttt{could\ not\ find\ function\ "read\_csv"} you have not loaded the
\texttt{tidyverse}.

After importing, always check that the resulting table looks like you
expect. Click \texttt{dat} in the Environment pane, or run
\texttt{glimpse(dat)} in the console.

\begin{Shaded}
\begin{Highlighting}[]
\FunctionTok{glimpse}\NormalTok{(dat)}
\end{Highlighting}
\end{Shaded}

\begin{verbatim}
Rows: 10,507
Columns: 13
$ `Voyage ID`                                                              <dbl> ~
$ `Vessel name`                                                            <chr> ~
$ `Voyage itinerary imputed port where began (ptdepimp) place`             <chr> ~
$ `Voyage itinerary imputed principal place of slave purchase (mjbyptimp)` <chr> ~
$ `Total embarked`                                                         <dbl> ~
$ `Total disembarked`                                                      <dbl> ~
$ `Year of arrival at port of disembarkation`                              <dbl> ~
$ `Percent men`                                                            <dbl> ~
$ `Percent women`                                                          <dbl> ~
$ `Percent children`                                                       <dbl> ~
$ `Slaves died during middle passage`                                      <dbl> ~
$ `Mortality rate`                                                         <dbl> ~
$ `Slaves outcome label`                                                   <chr> ~
\end{verbatim}

\begin{itemize}
\tightlist
\item
  The dataset has \_\_\_\_\_ rows of data and \_\_ columns.
\end{itemize}

\section{House-keeping}\label{house-keeping}

The column names are informative but long, and several contain spaces,
capital letters and brackets. That makes them awkward to type and to
refer to in code. As a rule, variable names should avoid spaces and
special characters. Pick a naming convention (\texttt{my.var},
\texttt{myVar}, \texttt{my\_var}) and stick with it. This tutorial uses
underscores and no capitals, a style often called snake case.

The \texttt{clean\_names()} function from the \texttt{janitor} package
does the mechanical part of this for us. It converts every column name
to lower case and replaces spaces and brackets with underscores, so we
no longer need back-ticks to refer to a column.

\begin{Shaded}
\begin{Highlighting}[]
\NormalTok{dat\_clean }\OtherTok{\textless{}{-}} \FunctionTok{clean\_names}\NormalTok{(dat)}

\FunctionTok{names}\NormalTok{(dat\_clean)}
\end{Highlighting}
\end{Shaded}

\begin{verbatim}
 [1] "voyage_id"                                                           
 [2] "vessel_name"                                                         
 [3] "voyage_itinerary_imputed_port_where_began_ptdepimp_place"            
 [4] "voyage_itinerary_imputed_principal_place_of_slave_purchase_mjbyptimp"
 [5] "total_embarked"                                                      
 [6] "total_disembarked"                                                   
 [7] "year_of_arrival_at_port_of_disembarkation"                           
 [8] "percent_men"                                                         
 [9] "percent_women"                                                       
[10] "percent_children"                                                    
[11] "slaves_died_during_middle_passage"                                   
[12] "mortality_rate"                                                      
[13] "slaves_outcome_label"                                                
\end{verbatim}

Every name is now consistent snake case. Two of them,
\texttt{total\_embarked} and \texttt{total\_disembarked}, are already
short enough to use as they are. The rest are still long, because
\texttt{clean\_names()} standardises names but cannot know which ones we
would like to be shorter.

For that we use \texttt{rename()}. The syntax is
\texttt{new\_name\ =\ old\_name}. Because \texttt{clean\_names()} has
already removed the spaces, no back-ticks are needed.

\begin{Shaded}
\begin{Highlighting}[]
\NormalTok{dat\_rename }\OtherTok{\textless{}{-}} \FunctionTok{rename}\NormalTok{(dat\_clean,}
                     \AttributeTok{id =}\NormalTok{ voyage\_id,}
                     \AttributeTok{vessel =}\NormalTok{ vessel\_name,}
                     \AttributeTok{start\_port =}\NormalTok{ voyage\_itinerary\_imputed\_port\_where\_began\_ptdepimp\_place,}
                     \AttributeTok{purchase\_place =}\NormalTok{ voyage\_itinerary\_imputed\_principal\_place\_of\_slave\_purchase\_mjbyptimp,}
                     \AttributeTok{arrival\_year =}\NormalTok{ year\_of\_arrival\_at\_port\_of\_disembarkation,}
                     \AttributeTok{men =}\NormalTok{ percent\_men,}
                     \AttributeTok{women =}\NormalTok{ percent\_women,}
                     \AttributeTok{children =}\NormalTok{ percent\_children,}
                     \AttributeTok{died =}\NormalTok{ slaves\_died\_during\_middle\_passage,}
                     \AttributeTok{mortality =}\NormalTok{ mortality\_rate,}
                     \AttributeTok{outcome =}\NormalTok{ slaves\_outcome\_label)}
\end{Highlighting}
\end{Shaded}

\begin{tcolorbox}[enhanced jigsaw, colframe=quarto-callout-tip-color-frame, breakable, left=2mm, leftrule=.75mm, toptitle=1mm, rightrule=.15mm, opacitybacktitle=0.6, arc=.35mm, colback=white, title=\textcolor{quarto-callout-tip-color}{\faLightbulb}\hspace{0.5em}{Tip}, coltitle=black, colbacktitle=quarto-callout-tip-color!10!white, toprule=.15mm, bottomtitle=1mm, opacityback=0, titlerule=0mm, bottomrule=.15mm]

When code starts to spread across many lines, pressing Enter after a
comma to break the line makes it much easier to read.

\end{tcolorbox}

\section{Checking variable types}\label{checking-variable-types}

After importing and renaming, check what R thinks each column is. Three
useful options are:

\begin{Shaded}
\begin{Highlighting}[]
\FunctionTok{str}\NormalTok{(dat\_rename)}
\FunctionTok{summary}\NormalTok{(dat\_rename)}
\FunctionTok{map\_chr}\NormalTok{(dat\_rename, class)}
\end{Highlighting}
\end{Shaded}

\begin{itemize}
\tightlist
\item
  How many of the variables are character (text) variables? \_
\item
  How many of the variables are numeric (double)? \_
\end{itemize}

Several are correct as-is, but two issues need attention. The
\texttt{id} column contains numbers, but those numbers are identifiers,
not quantities you would average. The character columns
(\texttt{vessel}, \texttt{start\_port}, \texttt{purchase\_place},
\texttt{outcome}) hold a fixed set of categories rather than free text.
Both kinds of column should be factors, which is how R represents
categorical data.

\section{Converting to factors}\label{converting-to-factors}

We use \texttt{mutate()} to overwrite each column with a factor version
of itself. \texttt{factor()} lets us set the display order with the
\texttt{levels} argument, which matters because R will otherwise sort
categories alphabetically in plots.

\begin{Shaded}
\begin{Highlighting}[]
\NormalTok{dat\_final }\OtherTok{\textless{}{-}} \FunctionTok{mutate}\NormalTok{(dat\_rename,}
                    \AttributeTok{id             =} \FunctionTok{factor}\NormalTok{(id),}
                    \AttributeTok{vessel         =} \FunctionTok{factor}\NormalTok{(vessel),}
                    \AttributeTok{purchase\_place =} \FunctionTok{factor}\NormalTok{(purchase\_place),}
                    \AttributeTok{outcome        =} \FunctionTok{factor}\NormalTok{(outcome),}
                    \AttributeTok{start\_port     =} \FunctionTok{factor}\NormalTok{(start\_port,}
                                            \AttributeTok{levels =} \FunctionTok{c}\NormalTok{(}\StringTok{"Liverpool"}\NormalTok{,}
                                                       \StringTok{"London"}\NormalTok{,}
                                                       \StringTok{"Bristol"}\NormalTok{,}
                                                       \FunctionTok{sort}\NormalTok{(}\FunctionTok{setdiff}\NormalTok{(}\FunctionTok{unique}\NormalTok{(start\_port),}
                                                                    \FunctionTok{c}\NormalTok{(}\StringTok{"Liverpool"}\NormalTok{, }\StringTok{"London"}\NormalTok{, }\StringTok{"Bristol"}\NormalTok{))))))}
\end{Highlighting}
\end{Shaded}

We have given \texttt{start\_port} an explicit order. Liverpool, London
and Bristol come first because we will focus on those three; the
remaining ports follow in alphabetical order. This means later plots
will show the three big British ports in a meaningful order, rather than
the default alphabetical Bristol, Liverpool, London.

\begin{tcolorbox}[enhanced jigsaw, colframe=quarto-callout-note-color-frame, breakable, left=2mm, leftrule=.75mm, toptitle=1mm, rightrule=.15mm, opacitybacktitle=0.6, arc=.35mm, colback=white, title=\textcolor{quarto-callout-note-color}{\faInfo}\hspace{0.5em}{A quicker way for batch conversions}, coltitle=black, colbacktitle=quarto-callout-note-color!10!white, toprule=.15mm, bottomtitle=1mm, opacityback=0, titlerule=0mm, bottomrule=.15mm]

If you do not need to control the levels for each column individually,
you can convert several columns at once with \texttt{across()}:

\begin{Shaded}
\begin{Highlighting}[]
\NormalTok{dat\_final }\OtherTok{\textless{}{-}} \FunctionTok{mutate}\NormalTok{(dat\_rename,}
                    \FunctionTok{across}\NormalTok{(}\FunctionTok{c}\NormalTok{(id, vessel, purchase\_place, outcome, start\_port),}
\NormalTok{                           as.factor))}
\end{Highlighting}
\end{Shaded}

This is more compact but loses the explicit ordering of
\texttt{start\_port}. See the
\href{https://dplyr.tidyverse.org/reference/across.html}{\texttt{dplyr}
documentation} for more.

\end{tcolorbox}

Run \texttt{summary()} again on \texttt{dat\_final} and look at how the
factor variables are now summarised by their counts per level.

\begin{Shaded}
\begin{Highlighting}[]
\FunctionTok{summary}\NormalTok{(dat\_final)}
\end{Highlighting}
\end{Shaded}

\section{Activities}\label{activities}

\begin{tcolorbox}[enhanced jigsaw, colframe=quarto-callout-tip-color-frame, breakable, left=2mm, leftrule=.75mm, toptitle=1mm, rightrule=.15mm, opacitybacktitle=0.6, arc=.35mm, colback=white, title=\textcolor{quarto-callout-tip-color}{\faLightbulb}\hspace{0.5em}{Activity 1}, coltitle=black, colbacktitle=quarto-callout-tip-color!10!white, toprule=.15mm, bottomtitle=1mm, opacityback=0, titlerule=0mm, bottomrule=.15mm]

Run \texttt{glimpse(dat\_final)} and identify which columns are now
factors. How many factor columns are there? \_

\end{tcolorbox}

\begin{tcolorbox}[enhanced jigsaw, colframe=quarto-callout-tip-color-frame, breakable, left=2mm, leftrule=.75mm, toptitle=1mm, rightrule=.15mm, opacitybacktitle=0.6, arc=.35mm, colback=white, title=\textcolor{quarto-callout-tip-color}{\faLightbulb}\hspace{0.5em}{Activity 2}, coltitle=black, colbacktitle=quarto-callout-tip-color!10!white, toprule=.15mm, bottomtitle=1mm, opacityback=0, titlerule=0mm, bottomrule=.15mm]

The \texttt{mortality} column is a proportion. Use
\texttt{summary(dat\_final\$mortality)} to find its mean. To two decimal
places, what is it? \_\_\_\_

\end{tcolorbox}

\begin{tcolorbox}[enhanced jigsaw, colframe=quarto-callout-note-color-frame, breakable, left=2mm, leftrule=.75mm, toptitle=1mm, rightrule=.15mm, opacitybacktitle=0.6, arc=.35mm, colback=white, title=\textcolor{quarto-callout-note-color}{\faInfo}\hspace{0.5em}{Solution to Activity 2}, coltitle=black, colbacktitle=quarto-callout-note-color!10!white, toprule=.15mm, bottomtitle=1mm, opacityback=0, titlerule=0mm, bottomrule=.15mm]

\begin{Shaded}
\begin{Highlighting}[]
\FunctionTok{summary}\NormalTok{(dat\_final}\SpecialCharTok{$}\NormalTok{mortality)}
\end{Highlighting}
\end{Shaded}

The mean mortality rate across the 10,507 voyages in this subset is
about 0.18 (18\%).

\end{tcolorbox}

\bookmarksetup{startatroot}

\chapter{Basic plots}\label{sec-basic-plots}

\section{Intended learning outcomes}\label{intended-learning-outcomes-2}

By the end of this chapter you will be able to:

\begin{itemize}
\tightlist
\item
  Explain what is meant by a ``grammar of graphics'' and a ``layer'' in
  \texttt{ggplot2}
\item
  Build a scatterplot using \texttt{geom\_point()}
\item
  Add a line of best fit using \texttt{geom\_smooth()}
\item
  Plot the distribution of a single continuous variable with
  \texttt{geom\_histogram()}
\item
  Plot counts with \texttt{geom\_bar()} and pre-computed values with
  \texttt{geom\_col()}
\item
  Plot grouped distributions with \texttt{geom\_boxplot()} and
  \texttt{geom\_violin()}
\end{itemize}

\section{Functions used}\label{functions-used-1}

\begin{itemize}
\tightlist
\item
  ggplot2: \texttt{ggplot()}, \texttt{aes()}, \texttt{geom\_point()},
  \texttt{geom\_smooth()}, \texttt{geom\_histogram()},
  \texttt{geom\_bar()}, \texttt{geom\_col()}, \texttt{geom\_boxplot()},
  \texttt{geom\_violin()}
\item
  dplyr: \texttt{filter()}, \texttt{group\_by()}, \texttt{count()},
  \texttt{arrange()}, \texttt{summarise()}
\end{itemize}

\section{The grammar of graphics}\label{the-grammar-of-graphics}

\texttt{ggplot2} is so powerful because it builds plots in layers. The
first layer is always the base of the plot, set up with
\texttt{ggplot()}, which specifies the data and the aesthetic mappings
(which variables go on which axis, or are mapped to colour or shape).
Subsequent layers add geoms (geometric objects such as points, bars, or
boxes), scales, themes and so on. Layers are added with \texttt{+}.

It is helpful to think of any plot as having several elements that sit
semi-transparently on top of each other. A scatterplot with a line of
best fit, for example, is the base layer, then the points layer, then
the smoother layer.

\section{Scatterplots}\label{scatterplots}

A scatterplot is appropriate when both axes are continuous. The base
layer below maps \texttt{arrival\_year} to the x-axis and
\texttt{total\_embarked} to the y-axis.

\begin{Shaded}
\begin{Highlighting}[]
\FunctionTok{ggplot}\NormalTok{(}\AttributeTok{data =}\NormalTok{ dat\_final, }\FunctionTok{aes}\NormalTok{(}\AttributeTok{x =}\NormalTok{ arrival\_year, }\AttributeTok{y =}\NormalTok{ total\_embarked))}
\end{Highlighting}
\end{Shaded}

Run that line on its own. You get an empty plot with the axes drawn but
no data. That is because we have not added a geom. Given the variable
types, what type of plot is likely to be built on this base layer?

\begin{itemize}
\tightlist
\item
  \begin{enumerate}
  \def\labelenumi{(\Alph{enumi})}
  \tightlist
  \item
    Bar chart\\
  \end{enumerate}
\item
  \begin{enumerate}
  \def\labelenumi{(\Alph{enumi})}
  \setcounter{enumi}{1}
  \tightlist
  \item
    Scatterplot\\
  \end{enumerate}
\item
  \begin{enumerate}
  \def\labelenumi{(\Alph{enumi})}
  \setcounter{enumi}{2}
  \tightlist
  \item
    Histogram
  \end{enumerate}
\end{itemize}

We add the points with \texttt{geom\_point()}.

\begin{Shaded}
\begin{Highlighting}[]
\FunctionTok{ggplot}\NormalTok{(dat\_final, }\FunctionTok{aes}\NormalTok{(arrival\_year, total\_embarked)) }\SpecialCharTok{+}
  \FunctionTok{geom\_point}\NormalTok{()}
\end{Highlighting}
\end{Shaded}

\begin{figure}[H]

\centering{

\includegraphics[width=1\linewidth,height=\textheight,keepaspectratio]{02-basic-plots_files/figure-pdf/fig-scatter-points-1.pdf}

}

\caption{\label{fig-scatter-points}Number embarked on each voyage,
plotted against the year the voyage arrived in port.}

\end{figure}%

\begin{tcolorbox}[enhanced jigsaw, colframe=quarto-callout-note-color-frame, breakable, left=2mm, leftrule=.75mm, toptitle=1mm, rightrule=.15mm, opacitybacktitle=0.6, arc=.35mm, colback=white, title=\textcolor{quarto-callout-note-color}{\faInfo}\hspace{0.5em}{Note}, coltitle=black, colbacktitle=quarto-callout-note-color!10!white, toprule=.15mm, bottomtitle=1mm, opacityback=0, titlerule=0mm, bottomrule=.15mm]

You will see a warning saying
\texttt{Removed\ 422\ rows\ containing\ missing\ values\ or\ values\ outside\ the\ scale\ range}.
That means 422 voyages in the data have a missing value for either
\texttt{arrival\_year} or \texttt{total\_embarked}. \texttt{ggplot2}
quietly drops them. We will look at handling missing data directly in
Section~\ref{sec-wide-to-long}.

\end{tcolorbox}

There are so many points that it is hard to see the trend. That is not a
fault of the plot; it is a reflection of the scale of the trade.

\section{Adding a line of best fit}\label{adding-a-line-of-best-fit}

A smoother makes it easier to see the trend. \texttt{geom\_smooth()}
uses LOESS by default for samples smaller than 1000 observations and GAM
otherwise. For a linear trend, set \texttt{method\ =\ "lm"}.

\begin{Shaded}
\begin{Highlighting}[]
\NormalTok{loess\_plot }\OtherTok{\textless{}{-}} \FunctionTok{ggplot}\NormalTok{(dat\_final, }\FunctionTok{aes}\NormalTok{(arrival\_year, total\_embarked)) }\SpecialCharTok{+}
  \FunctionTok{geom\_point}\NormalTok{(}\AttributeTok{alpha =} \FloatTok{0.3}\NormalTok{) }\SpecialCharTok{+}
  \FunctionTok{geom\_smooth}\NormalTok{()}

\NormalTok{linear\_plot }\OtherTok{\textless{}{-}} \FunctionTok{ggplot}\NormalTok{(dat\_final, }\FunctionTok{aes}\NormalTok{(arrival\_year, total\_embarked)) }\SpecialCharTok{+}
  \FunctionTok{geom\_point}\NormalTok{(}\AttributeTok{alpha =} \FloatTok{0.3}\NormalTok{) }\SpecialCharTok{+}
  \FunctionTok{geom\_smooth}\NormalTok{(}\AttributeTok{method =} \StringTok{"lm"}\NormalTok{, }\AttributeTok{formula =}\NormalTok{ y }\SpecialCharTok{\textasciitilde{}}\NormalTok{ x)}

\NormalTok{loess\_plot }\SpecialCharTok{+}\NormalTok{ linear\_plot}
\end{Highlighting}
\end{Shaded}

\begin{figure}[H]

\centering{

\includegraphics[width=1\linewidth,height=\textheight,keepaspectratio]{02-basic-plots_files/figure-pdf/fig-scatter-smooth-1.pdf}

}

\caption{\label{fig-scatter-smooth}Scatterplot with two smoothers, LOESS
(left) and linear (right).}

\end{figure}%

The combination of \texttt{loess\_plot\ +\ linear\_plot} is from the
\texttt{patchwork} package, which we cover properly in
Section~\ref{sec-patchwork}.

\section{Layer order matters}\label{layer-order-matters}

If you add \texttt{geom\_smooth()} first and \texttt{geom\_point()}
second, the points are drawn on top of the smoother.

\begin{Shaded}
\begin{Highlighting}[]
\FunctionTok{ggplot}\NormalTok{(dat\_final, }\FunctionTok{aes}\NormalTok{(arrival\_year, total\_embarked)) }\SpecialCharTok{+}
  \FunctionTok{geom\_smooth}\NormalTok{() }\SpecialCharTok{+}
  \FunctionTok{geom\_point}\NormalTok{()}
\end{Highlighting}
\end{Shaded}

\begin{figure}[H]

\centering{

\includegraphics[width=1\linewidth,height=\textheight,keepaspectratio]{02-basic-plots_files/figure-pdf/fig-scatter-wrong-order-1.pdf}

}

\caption{\label{fig-scatter-wrong-order}Layers in the wrong order:
points drawn on top of the smoother.}

\end{figure}%

\section{Histograms}\label{histograms}

A histogram shows the distribution of a single continuous variable. The
x-axis is divided into bins and the y-axis shows how many observations
fall into each bin.

\begin{Shaded}
\begin{Highlighting}[]
\FunctionTok{ggplot}\NormalTok{(dat\_final, }\FunctionTok{aes}\NormalTok{(}\AttributeTok{x =}\NormalTok{ total\_embarked)) }\SpecialCharTok{+}
  \FunctionTok{geom\_histogram}\NormalTok{()}
\end{Highlighting}
\end{Shaded}

\begin{figure}[H]

\centering{

\includegraphics[width=1\linewidth,height=\textheight,keepaspectratio]{02-basic-plots_files/figure-pdf/fig-histogram-default-1.pdf}

}

\caption{\label{fig-histogram-default}Default histogram of total
embarked per voyage.}

\end{figure}%

The default uses 30 bins and you will see the message
\texttt{stat\_bin()\ using\ bins\ =\ 30.\ Pick\ better\ value\ with\ binwidth}.
It is good practice to set \texttt{binwidth} (the size of each bin)
explicitly, along with \texttt{boundary\ =\ 0} so the bins start at
sensible round numbers.

\begin{Shaded}
\begin{Highlighting}[]
\FunctionTok{ggplot}\NormalTok{(dat\_final, }\FunctionTok{aes}\NormalTok{(}\AttributeTok{x =}\NormalTok{ total\_embarked)) }\SpecialCharTok{+}
  \FunctionTok{geom\_histogram}\NormalTok{(}\AttributeTok{binwidth =} \DecValTok{50}\NormalTok{, }\AttributeTok{boundary =} \DecValTok{0}\NormalTok{,}
                 \AttributeTok{fill =} \StringTok{"white"}\NormalTok{, }\AttributeTok{colour =} \StringTok{"black"}\NormalTok{)}
\end{Highlighting}
\end{Shaded}

\begin{figure}[H]

\centering{

\includegraphics[width=1\linewidth,height=\textheight,keepaspectratio]{02-basic-plots_files/figure-pdf/fig-histogram-binwidth-1.pdf}

}

\caption{\label{fig-histogram-binwidth}Histogram with explicit binwidth
of 50 starting at 0.}

\end{figure}%

\texttt{fill} sets the inside colour of each bar; \texttt{colour} sets
the outline.

\section{Bar charts of counts}\label{bar-charts-of-counts}

\texttt{geom\_bar()} counts the number of observations in each level of
its x-aesthetic. You do not specify a y-axis because the y-axis is
always a count.

If we try to plot \texttt{start\_port} directly, the result is unusable.

\begin{Shaded}
\begin{Highlighting}[]
\FunctionTok{ggplot}\NormalTok{(dat\_final, }\FunctionTok{aes}\NormalTok{(}\AttributeTok{x =}\NormalTok{ start\_port)) }\SpecialCharTok{+}
  \FunctionTok{geom\_bar}\NormalTok{()}
\end{Highlighting}
\end{Shaded}

There are too many ports to fit on one axis. We can confirm this with a
quick count.

\begin{Shaded}
\begin{Highlighting}[]
\NormalTok{port\_counts }\OtherTok{\textless{}{-}}\NormalTok{ dat\_final }\SpecialCharTok{|\textgreater{}}
  \FunctionTok{group\_by}\NormalTok{(start\_port) }\SpecialCharTok{|\textgreater{}}
  \FunctionTok{count}\NormalTok{() }\SpecialCharTok{|\textgreater{}}
  \FunctionTok{arrange}\NormalTok{(}\FunctionTok{desc}\NormalTok{(n))}

\NormalTok{port\_counts}
\end{Highlighting}
\end{Shaded}

\begin{verbatim}
# A tibble: 34 x 2
# Groups:   start_port [34]
   start_port                          n
   <fct>                           <int>
 1 Liverpool                        4973
 2 London                           3126
 3 Bristol                          2083
 4 Lancaster                         125
 5 Whitehaven                         64
 6 Plymouth                           32
 7 Great Britain, port unspecified    15
 8 Greenock                           11
 9 Poole                               9
10 Chester                             7
# i 24 more rows
\end{verbatim}

Liverpool, London and Bristol dominate. We will filter to those three
for the next several plots.

\begin{Shaded}
\begin{Highlighting}[]
\NormalTok{dat\_filter }\OtherTok{\textless{}{-}}\NormalTok{ dat\_final }\SpecialCharTok{|\textgreater{}}
  \FunctionTok{filter}\NormalTok{(start\_port }\SpecialCharTok{\%in\%} \FunctionTok{c}\NormalTok{(}\StringTok{"Liverpool"}\NormalTok{, }\StringTok{"London"}\NormalTok{, }\StringTok{"Bristol"}\NormalTok{))}
\end{Highlighting}
\end{Shaded}

The \texttt{\%in\%} operator reads as ``is one of''.
\texttt{dat\_filter} keeps only the rows where \texttt{start\_port}
matches one of the three values.

\begin{Shaded}
\begin{Highlighting}[]
\FunctionTok{ggplot}\NormalTok{(dat\_filter, }\FunctionTok{aes}\NormalTok{(}\AttributeTok{x =}\NormalTok{ start\_port, }\AttributeTok{fill =}\NormalTok{ start\_port)) }\SpecialCharTok{+}
  \FunctionTok{geom\_bar}\NormalTok{(}\AttributeTok{show.legend =} \ConstantTok{FALSE}\NormalTok{)}
\end{Highlighting}
\end{Shaded}

\begin{figure}[H]

\centering{

\includegraphics[width=1\linewidth,height=\textheight,keepaspectratio]{02-basic-plots_files/figure-pdf/fig-bar-ports-1.pdf}

}

\caption{\label{fig-bar-ports}Number of voyages originating in each of
the three main British ports.}

\end{figure}%

\begin{tcolorbox}[enhanced jigsaw, colframe=quarto-callout-warning-color-frame, breakable, left=2mm, leftrule=.75mm, toptitle=1mm, rightrule=.15mm, opacitybacktitle=0.6, arc=.35mm, colback=white, title=\textcolor{quarto-callout-warning-color}{\faExclamationTriangle}\hspace{0.5em}{Warning}, coltitle=black, colbacktitle=quarto-callout-warning-color!10!white, toprule=.15mm, bottomtitle=1mm, opacityback=0, titlerule=0mm, bottomrule=.15mm]

This bar chart shows the number of \emph{voyages}, not the number of
people. Each Liverpool voyage trafficked, on average, hundreds of
people, so the human cost is far larger than the bar heights suggest. We
will plot the human totals shortly.

\end{tcolorbox}

\section{\texorpdfstring{Bar charts of pre-computed values
(\texttt{geom\_col()})}{Bar charts of pre-computed values (geom\_col())}}\label{bar-charts-of-pre-computed-values-geom_col}

When you already have the y-values (because you computed them yourself),
use \texttt{geom\_col()} instead of \texttt{geom\_bar()}. The default
statistic for \texttt{geom\_col()} is \texttt{"identity"}: it uses the
numbers you supplied as-is.

\begin{Shaded}
\begin{Highlighting}[]
\NormalTok{port\_summary }\OtherTok{\textless{}{-}}\NormalTok{ dat\_filter }\SpecialCharTok{|\textgreater{}}
  \FunctionTok{group\_by}\NormalTok{(start\_port) }\SpecialCharTok{|\textgreater{}}
  \FunctionTok{summarise}\NormalTok{(}\AttributeTok{embarked =} \FunctionTok{sum}\NormalTok{(total\_embarked, }\AttributeTok{na.rm =} \ConstantTok{TRUE}\NormalTok{),}
            \AttributeTok{.groups =} \StringTok{"drop"}\NormalTok{)}

\NormalTok{port\_summary}
\end{Highlighting}
\end{Shaded}

\begin{verbatim}
# A tibble: 3 x 2
  start_port embarked
  <fct>         <dbl>
1 Liverpool   1337530
2 London       832944
3 Bristol      563796
\end{verbatim}

\begin{itemize}
\item
  \texttt{na.rm\ =\ TRUE} tells \texttt{sum()} to ignore missing values
  when computing the total. Without it, any single missing value would
  produce \texttt{NA} for the whole group.
\item
  \texttt{.groups\ =\ "drop"} removes the grouping after summarising,
  which is the modern \texttt{dplyr} recommendation; without it you get
  a ``regrouping'' message.
\item
  How many enslaved people embarked on ships originating in Liverpool?
\item
  \begin{enumerate}
  \def\labelenumi{(\Alph{enumi})}
  \tightlist
  \item
    \textasciitilde832 thousand\\
  \end{enumerate}
\item
  \begin{enumerate}
  \def\labelenumi{(\Alph{enumi})}
  \setcounter{enumi}{1}
  \tightlist
  \item
    \textasciitilde1.3 million\\
  \end{enumerate}
\item
  \begin{enumerate}
  \def\labelenumi{(\Alph{enumi})}
  \setcounter{enumi}{2}
  \tightlist
  \item
    56 thousand
  \end{enumerate}
\end{itemize}

\begin{Shaded}
\begin{Highlighting}[]
\FunctionTok{ggplot}\NormalTok{(port\_summary, }\FunctionTok{aes}\NormalTok{(}\AttributeTok{x =}\NormalTok{ start\_port, }\AttributeTok{y =}\NormalTok{ embarked, }\AttributeTok{fill =}\NormalTok{ start\_port)) }\SpecialCharTok{+}
  \FunctionTok{geom\_col}\NormalTok{(}\AttributeTok{show.legend =} \ConstantTok{FALSE}\NormalTok{)}
\end{Highlighting}
\end{Shaded}

\begin{figure}[H]

\centering{

\includegraphics[width=1\linewidth,height=\textheight,keepaspectratio]{02-basic-plots_files/figure-pdf/fig-col-totals-1.pdf}

}

\caption{\label{fig-col-totals}Total number of enslaved people embarked
on ships originating in each port.}

\end{figure}%

\section{Boxplots}\label{boxplots}

A boxplot shows the median, interquartile range and outliers of a
continuous variable, optionally broken down by a categorical variable on
the x-axis.

\begin{Shaded}
\begin{Highlighting}[]
\FunctionTok{ggplot}\NormalTok{(dat\_filter, }\FunctionTok{aes}\NormalTok{(}\AttributeTok{x =}\NormalTok{ start\_port, }\AttributeTok{y =}\NormalTok{ total\_embarked)) }\SpecialCharTok{+}
  \FunctionTok{geom\_boxplot}\NormalTok{()}
\end{Highlighting}
\end{Shaded}

\begin{figure}[H]

\centering{

\includegraphics[width=1\linewidth,height=\textheight,keepaspectratio]{02-basic-plots_files/figure-pdf/fig-boxplot-1.pdf}

}

\caption{\label{fig-boxplot}Distribution of the number embarked per
voyage, by port of origin.}

\end{figure}%

The median number of people embarked per voyage was just over 250, but
in many cases the number was much higher, particularly for ships
originating in Liverpool.

\section{Violin-boxplots}\label{violin-boxplots}

A violin plot shows the full distribution as a mirrored density.
Combined with a thin boxplot inside it, the result is much more
informative than a boxplot alone.

\begin{Shaded}
\begin{Highlighting}[]
\FunctionTok{ggplot}\NormalTok{(dat\_filter, }\FunctionTok{aes}\NormalTok{(}\AttributeTok{x =}\NormalTok{ start\_port, }\AttributeTok{y =}\NormalTok{ total\_embarked)) }\SpecialCharTok{+}
  \FunctionTok{geom\_violin}\NormalTok{(}\AttributeTok{trim =} \ConstantTok{FALSE}\NormalTok{) }\SpecialCharTok{+}
  \FunctionTok{geom\_boxplot}\NormalTok{(}\AttributeTok{width =} \FloatTok{0.2}\NormalTok{)}
\end{Highlighting}
\end{Shaded}

\begin{figure}[H]

\centering{

\includegraphics[width=1\linewidth,height=\textheight,keepaspectratio]{02-basic-plots_files/figure-pdf/fig-violin-boxplot-1.pdf}

}

\caption{\label{fig-violin-boxplot}Violin-boxplots of the number
embarked per voyage.}

\end{figure}%

\begin{tcolorbox}[enhanced jigsaw, colframe=quarto-callout-tip-color-frame, breakable, left=2mm, leftrule=.75mm, toptitle=1mm, rightrule=.15mm, opacitybacktitle=0.6, arc=.35mm, colback=white, title=\textcolor{quarto-callout-tip-color}{\faLightbulb}\hspace{0.5em}{Tip}, coltitle=black, colbacktitle=quarto-callout-tip-color!10!white, toprule=.15mm, bottomtitle=1mm, opacityback=0, titlerule=0mm, bottomrule=.15mm]

The order of the layers matters. If you draw the boxplot first and the
violin on top, the boxplot is hidden. Try changing
\texttt{trim\ =\ FALSE} to \texttt{TRUE} and \texttt{width\ =\ 0.2} to
\texttt{0.5} to see what each argument does.

\end{tcolorbox}

\section{Activities}\label{activities-1}

\begin{tcolorbox}[enhanced jigsaw, colframe=quarto-callout-tip-color-frame, breakable, left=2mm, leftrule=.75mm, toptitle=1mm, rightrule=.15mm, opacitybacktitle=0.6, arc=.35mm, colback=white, title=\textcolor{quarto-callout-tip-color}{\faLightbulb}\hspace{0.5em}{Activity 1}, coltitle=black, colbacktitle=quarto-callout-tip-color!10!white, toprule=.15mm, bottomtitle=1mm, opacityback=0, titlerule=0mm, bottomrule=.15mm]

Modify the histogram so the binwidth is 25 instead of 50. Does the peak
shift?

\end{tcolorbox}

\begin{tcolorbox}[enhanced jigsaw, colframe=quarto-callout-note-color-frame, breakable, left=2mm, leftrule=.75mm, toptitle=1mm, rightrule=.15mm, opacitybacktitle=0.6, arc=.35mm, colback=white, title=\textcolor{quarto-callout-note-color}{\faInfo}\hspace{0.5em}{Solution}, coltitle=black, colbacktitle=quarto-callout-note-color!10!white, toprule=.15mm, bottomtitle=1mm, opacityback=0, titlerule=0mm, bottomrule=.15mm]

\begin{Shaded}
\begin{Highlighting}[]
\FunctionTok{ggplot}\NormalTok{(dat\_final, }\FunctionTok{aes}\NormalTok{(}\AttributeTok{x =}\NormalTok{ total\_embarked)) }\SpecialCharTok{+}
  \FunctionTok{geom\_histogram}\NormalTok{(}\AttributeTok{binwidth =} \DecValTok{25}\NormalTok{, }\AttributeTok{boundary =} \DecValTok{0}\NormalTok{,}
                 \AttributeTok{fill =} \StringTok{"white"}\NormalTok{, }\AttributeTok{colour =} \StringTok{"black"}\NormalTok{)}
\end{Highlighting}
\end{Shaded}

The peak stays in roughly the same place (200-300) but the shape is more
detailed. Smaller bins reveal structure that larger bins smooth over, at
the cost of more visual noise.

\end{tcolorbox}

\begin{tcolorbox}[enhanced jigsaw, colframe=quarto-callout-tip-color-frame, breakable, left=2mm, leftrule=.75mm, toptitle=1mm, rightrule=.15mm, opacitybacktitle=0.6, arc=.35mm, colback=white, title=\textcolor{quarto-callout-tip-color}{\faLightbulb}\hspace{0.5em}{Activity 2}, coltitle=black, colbacktitle=quarto-callout-tip-color!10!white, toprule=.15mm, bottomtitle=1mm, opacityback=0, titlerule=0mm, bottomrule=.15mm]

Build a boxplot of \texttt{mortality} by \texttt{start\_port} (still
using \texttt{dat\_filter}).

\end{tcolorbox}

\begin{tcolorbox}[enhanced jigsaw, colframe=quarto-callout-note-color-frame, breakable, left=2mm, leftrule=.75mm, toptitle=1mm, rightrule=.15mm, opacitybacktitle=0.6, arc=.35mm, colback=white, title=\textcolor{quarto-callout-note-color}{\faInfo}\hspace{0.5em}{Solution}, coltitle=black, colbacktitle=quarto-callout-note-color!10!white, toprule=.15mm, bottomtitle=1mm, opacityback=0, titlerule=0mm, bottomrule=.15mm]

\begin{Shaded}
\begin{Highlighting}[]
\FunctionTok{ggplot}\NormalTok{(dat\_filter, }\FunctionTok{aes}\NormalTok{(}\AttributeTok{x =}\NormalTok{ start\_port, }\AttributeTok{y =}\NormalTok{ mortality)) }\SpecialCharTok{+}
  \FunctionTok{geom\_boxplot}\NormalTok{()}
\end{Highlighting}
\end{Shaded}

\begin{figure}[H]

\centering{

\includegraphics[width=1\linewidth,height=\textheight,keepaspectratio]{02-basic-plots_files/figure-pdf/fig-mortality-boxplot-1.pdf}

}

\caption{\label{fig-mortality-boxplot}Mortality rate per voyage by port
of origin.}

\end{figure}%

\end{tcolorbox}

\bookmarksetup{startatroot}

\chapter{Customising plots}\label{sec-customising}

\section{Intended learning outcomes}\label{intended-learning-outcomes-3}

By the end of this chapter you will be able to:

\begin{itemize}
\tightlist
\item
  Use \texttt{scale\_*} functions to control axis labels, breaks and
  limits
\item
  Distinguish \texttt{coord\_cartesian()} from \texttt{scale\_*} limits
\item
  Apply built-in themes and set a global default with
  \texttt{theme\_set()}
\item
  Add titles, subtitles and captions with \texttt{labs()}
\item
  Use the colour-blind-friendly \texttt{viridis} palettes
\end{itemize}

\section{Functions used}\label{functions-used-2}

\begin{itemize}
\tightlist
\item
  ggplot2: \texttt{scale\_x\_discrete()},
  \texttt{scale\_x\_continuous()}, \texttt{scale\_y\_continuous()},
  \texttt{coord\_cartesian()}, \texttt{scale\_fill\_viridis\_d()},
  \texttt{scale\_colour\_viridis\_d()}, \texttt{labs()},
  \texttt{theme\_minimal()}, \texttt{theme\_bw()},
  \texttt{theme\_classic()}, \texttt{theme\_set()}
\item
  scales: \texttt{label\_comma()} (replacement for
  \texttt{scales::comma})
\end{itemize}

\section{Scientific notation and number
formatting}\label{scientific-notation-and-number-formatting}

The column chart from Chapter~\ref{sec-basic-plots} uses scientific
notation on the y-axis because the numbers are so large.

\begin{Shaded}
\begin{Highlighting}[]
\FunctionTok{ggplot}\NormalTok{(port\_summary, }\FunctionTok{aes}\NormalTok{(}\AttributeTok{x =}\NormalTok{ start\_port, }\AttributeTok{y =}\NormalTok{ embarked, }\AttributeTok{fill =}\NormalTok{ start\_port)) }\SpecialCharTok{+}
  \FunctionTok{geom\_col}\NormalTok{(}\AttributeTok{show.legend =} \ConstantTok{FALSE}\NormalTok{) }\SpecialCharTok{+}
  \FunctionTok{scale\_fill\_viridis\_d}\NormalTok{(}\AttributeTok{option =} \StringTok{"E"}\NormalTok{)}
\end{Highlighting}
\end{Shaded}

\begin{figure}[H]

\centering{

\includegraphics[width=1\linewidth,height=\textheight,keepaspectratio]{03-customising_files/figure-pdf/fig-col-scientific-1.pdf}

}

\caption{\label{fig-col-scientific}Total embarked, with the default
y-axis in scientific notation.}

\end{figure}%

We can override that with \texttt{scale\_y\_continuous()} and a
label-formatting function from the \texttt{scales} package.

\begin{Shaded}
\begin{Highlighting}[]
\FunctionTok{ggplot}\NormalTok{(port\_summary, }\FunctionTok{aes}\NormalTok{(}\AttributeTok{x =}\NormalTok{ start\_port, }\AttributeTok{y =}\NormalTok{ embarked, }\AttributeTok{fill =}\NormalTok{ start\_port)) }\SpecialCharTok{+}
  \FunctionTok{geom\_col}\NormalTok{(}\AttributeTok{show.legend =} \ConstantTok{FALSE}\NormalTok{) }\SpecialCharTok{+}
  \FunctionTok{scale\_fill\_viridis\_d}\NormalTok{(}\AttributeTok{option =} \StringTok{"E"}\NormalTok{) }\SpecialCharTok{+}
  \FunctionTok{scale\_y\_continuous}\NormalTok{(}\AttributeTok{labels =}\NormalTok{ scales}\SpecialCharTok{::}\FunctionTok{label\_comma}\NormalTok{())}
\end{Highlighting}
\end{Shaded}

\begin{figure}[H]

\centering{

\includegraphics[width=1\linewidth,height=\textheight,keepaspectratio]{03-customising_files/figure-pdf/fig-col-comma-1.pdf}

}

\caption{\label{fig-col-comma}Total embarked, with the y-axis formatted
with thousands separators.}

\end{figure}%

\begin{tcolorbox}[enhanced jigsaw, colframe=quarto-callout-note-color-frame, breakable, left=2mm, leftrule=.75mm, toptitle=1mm, rightrule=.15mm, opacitybacktitle=0.6, arc=.35mm, colback=white, title=\textcolor{quarto-callout-note-color}{\faInfo}\hspace{0.5em}{Note}, coltitle=black, colbacktitle=quarto-callout-note-color!10!white, toprule=.15mm, bottomtitle=1mm, opacityback=0, titlerule=0mm, bottomrule=.15mm]

\texttt{scales::label\_comma()} is the modern factory function. The
older \texttt{scales::comma} (no parentheses) still works but is being
phased out.

\end{tcolorbox}

\section{Axis labels and titles}\label{axis-labels-and-titles}

\texttt{labs()} is the most concise way to set axis labels, the plot
title, subtitle and caption.

\begin{Shaded}
\begin{Highlighting}[]
\FunctionTok{ggplot}\NormalTok{(port\_summary, }\FunctionTok{aes}\NormalTok{(}\AttributeTok{x =}\NormalTok{ start\_port, }\AttributeTok{y =}\NormalTok{ embarked, }\AttributeTok{fill =}\NormalTok{ start\_port)) }\SpecialCharTok{+}
  \FunctionTok{geom\_col}\NormalTok{(}\AttributeTok{show.legend =} \ConstantTok{FALSE}\NormalTok{) }\SpecialCharTok{+}
  \FunctionTok{scale\_fill\_viridis\_d}\NormalTok{(}\AttributeTok{option =} \StringTok{"E"}\NormalTok{) }\SpecialCharTok{+}
  \FunctionTok{scale\_y\_continuous}\NormalTok{(}\AttributeTok{labels =}\NormalTok{ scales}\SpecialCharTok{::}\FunctionTok{label\_comma}\NormalTok{()) }\SpecialCharTok{+}
  \FunctionTok{labs}\NormalTok{(}\AttributeTok{x =} \StringTok{"Port of origin"}\NormalTok{,}
       \AttributeTok{y =} \StringTok{"Number embarked"}\NormalTok{,}
       \AttributeTok{title =} \StringTok{"Number embarked on ships originating in British ports"}\NormalTok{)}
\end{Highlighting}
\end{Shaded}

\begin{figure}[H]

\centering{

\includegraphics[width=1\linewidth,height=\textheight,keepaspectratio]{03-customising_files/figure-pdf/fig-col-labs-1.pdf}

}

\caption{\label{fig-col-labs}Adding axis labels and a title.}

\end{figure}%

\section{Themes}\label{themes}

\texttt{ggplot2} has several built-in themes. We have already set
\texttt{theme\_minimal()} as the global default in the project setup
(see \texttt{R/my\_setup.R}). You can override it for an individual plot
by adding a different theme as a layer.

\begin{Shaded}
\begin{Highlighting}[]
\NormalTok{p }\OtherTok{\textless{}{-}} \FunctionTok{ggplot}\NormalTok{(port\_summary, }\FunctionTok{aes}\NormalTok{(}\AttributeTok{x =}\NormalTok{ start\_port, }\AttributeTok{y =}\NormalTok{ embarked, }\AttributeTok{fill =}\NormalTok{ start\_port)) }\SpecialCharTok{+}
  \FunctionTok{geom\_col}\NormalTok{(}\AttributeTok{show.legend =} \ConstantTok{FALSE}\NormalTok{) }\SpecialCharTok{+}
  \FunctionTok{scale\_fill\_viridis\_d}\NormalTok{(}\AttributeTok{option =} \StringTok{"E"}\NormalTok{) }\SpecialCharTok{+}
  \FunctionTok{scale\_y\_continuous}\NormalTok{(}\AttributeTok{labels =}\NormalTok{ scales}\SpecialCharTok{::}\FunctionTok{label\_comma}\NormalTok{()) }\SpecialCharTok{+}
  \FunctionTok{labs}\NormalTok{(}\AttributeTok{x =} \ConstantTok{NULL}\NormalTok{, }\AttributeTok{y =} \ConstantTok{NULL}\NormalTok{)}

\NormalTok{(p }\SpecialCharTok{+} \FunctionTok{theme\_minimal}\NormalTok{() }\SpecialCharTok{+} \FunctionTok{ggtitle}\NormalTok{(}\StringTok{"theme\_minimal"}\NormalTok{)) }\SpecialCharTok{+}
\NormalTok{  (p }\SpecialCharTok{+} \FunctionTok{theme\_bw}\NormalTok{() }\SpecialCharTok{+} \FunctionTok{ggtitle}\NormalTok{(}\StringTok{"theme\_bw"}\NormalTok{)) }\SpecialCharTok{+}
\NormalTok{  (p }\SpecialCharTok{+} \FunctionTok{theme\_classic}\NormalTok{() }\SpecialCharTok{+} \FunctionTok{ggtitle}\NormalTok{(}\StringTok{"theme\_classic"}\NormalTok{))}
\end{Highlighting}
\end{Shaded}

\begin{figure}[H]

\centering{

\includegraphics[width=1\linewidth,height=\textheight,keepaspectratio]{03-customising_files/figure-pdf/fig-themes-1.pdf}

}

\caption{\label{fig-themes}The same plot in three different themes.}

\end{figure}%

\begin{tcolorbox}[enhanced jigsaw, colframe=quarto-callout-tip-color-frame, breakable, left=2mm, leftrule=.75mm, toptitle=1mm, rightrule=.15mm, opacitybacktitle=0.6, arc=.35mm, colback=white, title=\textcolor{quarto-callout-tip-color}{\faLightbulb}\hspace{0.5em}{Tip}, coltitle=black, colbacktitle=quarto-callout-tip-color!10!white, toprule=.15mm, bottomtitle=1mm, opacityback=0, titlerule=0mm, bottomrule=.15mm]

Type \texttt{theme\_} in the RStudio source pane and let auto-complete
show you the available options. To set a theme for every subsequent plot
in your script, run \texttt{theme\_set(theme\_bw())} near the top.

\end{tcolorbox}

\section{Discrete x-axis labels}\label{discrete-x-axis-labels}

\texttt{scale\_x\_discrete()} lets you control the tick labels on a
categorical axis without renaming the underlying data.

\begin{Shaded}
\begin{Highlighting}[]
\FunctionTok{ggplot}\NormalTok{(port\_summary, }\FunctionTok{aes}\NormalTok{(}\AttributeTok{x =}\NormalTok{ start\_port, }\AttributeTok{y =}\NormalTok{ embarked, }\AttributeTok{fill =}\NormalTok{ start\_port)) }\SpecialCharTok{+}
  \FunctionTok{geom\_col}\NormalTok{(}\AttributeTok{show.legend =} \ConstantTok{FALSE}\NormalTok{) }\SpecialCharTok{+}
  \FunctionTok{scale\_x\_discrete}\NormalTok{(}\AttributeTok{labels =} \FunctionTok{c}\NormalTok{(}\StringTok{"Liverpool"}\NormalTok{, }\StringTok{"Port of London"}\NormalTok{, }\StringTok{"Bristol"}\NormalTok{)) }\SpecialCharTok{+}
  \FunctionTok{scale\_fill\_viridis\_d}\NormalTok{(}\AttributeTok{option =} \StringTok{"E"}\NormalTok{) }\SpecialCharTok{+}
  \FunctionTok{scale\_y\_continuous}\NormalTok{(}\AttributeTok{labels =}\NormalTok{ scales}\SpecialCharTok{::}\FunctionTok{label\_comma}\NormalTok{())}
\end{Highlighting}
\end{Shaded}

\begin{figure}[H]

\centering{

\includegraphics[width=1\linewidth,height=\textheight,keepaspectratio]{03-customising_files/figure-pdf/fig-discrete-labels-1.pdf}

}

\caption{\label{fig-discrete-labels}Renaming x-axis ticks with
scale\_x\_discrete().}

\end{figure}%

\section{Axis limits the right way}\label{axis-limits-the-right-way}

There are two ways to restrict the visible range of an axis, and they
behave differently.

\texttt{coord\_cartesian(xlim\ =\ c(...),\ ylim\ =\ c(...))} zooms in on
the plot. Data outside the range are still used in calculations (such as
drawing a smoother) and are simply cropped from view.

\texttt{scale\_x\_continuous(limits\ =\ c(...))} filters the data. Any
rows outside the range are dropped before geoms or smoothers see them.
This can silently change a regression line, so it is usually the wrong
choice.

\begin{Shaded}
\begin{Highlighting}[]
\NormalTok{crop }\OtherTok{\textless{}{-}} \FunctionTok{ggplot}\NormalTok{(dat\_filter, }\FunctionTok{aes}\NormalTok{(arrival\_year, total\_embarked)) }\SpecialCharTok{+}
  \FunctionTok{geom\_point}\NormalTok{(}\AttributeTok{alpha =} \FloatTok{0.2}\NormalTok{) }\SpecialCharTok{+}
  \FunctionTok{geom\_smooth}\NormalTok{() }\SpecialCharTok{+}
  \FunctionTok{coord\_cartesian}\NormalTok{(}\AttributeTok{xlim =} \FunctionTok{c}\NormalTok{(}\DecValTok{1700}\NormalTok{, }\DecValTok{1815}\NormalTok{)) }\SpecialCharTok{+}
  \FunctionTok{labs}\NormalTok{(}\AttributeTok{title =} \StringTok{"coord\_cartesian (crop)"}\NormalTok{)}

\NormalTok{filt }\OtherTok{\textless{}{-}} \FunctionTok{ggplot}\NormalTok{(dat\_filter, }\FunctionTok{aes}\NormalTok{(arrival\_year, total\_embarked)) }\SpecialCharTok{+}
  \FunctionTok{geom\_point}\NormalTok{(}\AttributeTok{alpha =} \FloatTok{0.2}\NormalTok{) }\SpecialCharTok{+}
  \FunctionTok{geom\_smooth}\NormalTok{() }\SpecialCharTok{+}
  \FunctionTok{scale\_x\_continuous}\NormalTok{(}\AttributeTok{limits =} \FunctionTok{c}\NormalTok{(}\DecValTok{1700}\NormalTok{, }\DecValTok{1815}\NormalTok{)) }\SpecialCharTok{+}
  \FunctionTok{labs}\NormalTok{(}\AttributeTok{title =} \StringTok{"scale\_x\_continuous(limits=)"}\NormalTok{)}

\NormalTok{crop }\SpecialCharTok{+}\NormalTok{ filt}
\end{Highlighting}
\end{Shaded}

\begin{figure}[H]

\centering{

\includegraphics[width=1\linewidth,height=\textheight,keepaspectratio]{03-customising_files/figure-pdf/fig-axis-limits-1.pdf}

}

\caption{\label{fig-axis-limits}coord\_cartesian (left) crops the view;
scale limits (right) filter the data.}

\end{figure}%

\section{Viridis colour palettes}\label{viridis-colour-palettes}

The default \texttt{ggplot2} colours are not colour-blind friendly. The
\texttt{viridis} family of palettes is and is now included in
\texttt{ggplot2}. Use \texttt{scale\_fill\_viridis\_d()} for discrete
fill (a category mapped to fill), and \texttt{scale\_fill\_viridis\_c()}
for continuous fill (a number mapped to fill); equivalents exist for
\texttt{colour}.

The \texttt{option} argument picks one of five palettes:

\begin{itemize}
\tightlist
\item
  \texttt{"A"} magma
\item
  \texttt{"B"} inferno
\item
  \texttt{"C"} plasma
\item
  \texttt{"D"} viridis (the default)
\item
  \texttt{"E"} cividis (the most accessible)
\end{itemize}

\begin{Shaded}
\begin{Highlighting}[]
\NormalTok{make\_plot }\OtherTok{\textless{}{-}} \ControlFlowTok{function}\NormalTok{(opt) \{}
  \FunctionTok{ggplot}\NormalTok{(port\_summary, }\FunctionTok{aes}\NormalTok{(start\_port, embarked, }\AttributeTok{fill =}\NormalTok{ start\_port)) }\SpecialCharTok{+}
    \FunctionTok{geom\_col}\NormalTok{(}\AttributeTok{show.legend =} \ConstantTok{FALSE}\NormalTok{) }\SpecialCharTok{+}
    \FunctionTok{scale\_fill\_viridis\_d}\NormalTok{(}\AttributeTok{option =}\NormalTok{ opt) }\SpecialCharTok{+}
    \FunctionTok{labs}\NormalTok{(}\AttributeTok{title =} \FunctionTok{paste}\NormalTok{(}\StringTok{"option ="}\NormalTok{, opt), }\AttributeTok{x =} \ConstantTok{NULL}\NormalTok{, }\AttributeTok{y =} \ConstantTok{NULL}\NormalTok{)}
\NormalTok{\}}

\FunctionTok{make\_plot}\NormalTok{(}\StringTok{"A"}\NormalTok{) }\SpecialCharTok{+} \FunctionTok{make\_plot}\NormalTok{(}\StringTok{"B"}\NormalTok{) }\SpecialCharTok{+} \FunctionTok{make\_plot}\NormalTok{(}\StringTok{"C"}\NormalTok{) }\SpecialCharTok{+}
  \FunctionTok{make\_plot}\NormalTok{(}\StringTok{"D"}\NormalTok{) }\SpecialCharTok{+} \FunctionTok{make\_plot}\NormalTok{(}\StringTok{"E"}\NormalTok{) }\SpecialCharTok{+}
  \FunctionTok{plot\_layout}\NormalTok{(}\AttributeTok{nrow =} \DecValTok{2}\NormalTok{)}
\end{Highlighting}
\end{Shaded}

\begin{figure}[H]

\centering{

\includegraphics[width=1\linewidth,height=\textheight,keepaspectratio]{03-customising_files/figure-pdf/fig-viridis-options-1.pdf}

}

\caption{\label{fig-viridis-options}Five viridis palettes applied to the
same column chart.}

\end{figure}%

\section{Activities}\label{activities-2}

\begin{tcolorbox}[enhanced jigsaw, colframe=quarto-callout-tip-color-frame, breakable, left=2mm, leftrule=.75mm, toptitle=1mm, rightrule=.15mm, opacitybacktitle=0.6, arc=.35mm, colback=white, title=\textcolor{quarto-callout-tip-color}{\faLightbulb}\hspace{0.5em}{Activity 1}, coltitle=black, colbacktitle=quarto-callout-tip-color!10!white, toprule=.15mm, bottomtitle=1mm, opacityback=0, titlerule=0mm, bottomrule=.15mm]

Apply the \texttt{theme\_classic()} theme to the column chart of port
totals and edit the y-axis label to ``Total embarked (people)''.

\end{tcolorbox}

\begin{tcolorbox}[enhanced jigsaw, colframe=quarto-callout-note-color-frame, breakable, left=2mm, leftrule=.75mm, toptitle=1mm, rightrule=.15mm, opacitybacktitle=0.6, arc=.35mm, colback=white, title=\textcolor{quarto-callout-note-color}{\faInfo}\hspace{0.5em}{Solution}, coltitle=black, colbacktitle=quarto-callout-note-color!10!white, toprule=.15mm, bottomtitle=1mm, opacityback=0, titlerule=0mm, bottomrule=.15mm]

\begin{Shaded}
\begin{Highlighting}[]
\FunctionTok{ggplot}\NormalTok{(port\_summary, }\FunctionTok{aes}\NormalTok{(}\AttributeTok{x =}\NormalTok{ start\_port, }\AttributeTok{y =}\NormalTok{ embarked, }\AttributeTok{fill =}\NormalTok{ start\_port)) }\SpecialCharTok{+}
  \FunctionTok{geom\_col}\NormalTok{(}\AttributeTok{show.legend =} \ConstantTok{FALSE}\NormalTok{) }\SpecialCharTok{+}
  \FunctionTok{scale\_fill\_viridis\_d}\NormalTok{(}\AttributeTok{option =} \StringTok{"E"}\NormalTok{) }\SpecialCharTok{+}
  \FunctionTok{scale\_y\_continuous}\NormalTok{(}\AttributeTok{labels =}\NormalTok{ scales}\SpecialCharTok{::}\FunctionTok{label\_comma}\NormalTok{()) }\SpecialCharTok{+}
  \FunctionTok{labs}\NormalTok{(}\AttributeTok{x =} \StringTok{"Port of origin"}\NormalTok{, }\AttributeTok{y =} \StringTok{"Total embarked (people)"}\NormalTok{) }\SpecialCharTok{+}
  \FunctionTok{theme\_classic}\NormalTok{()}
\end{Highlighting}
\end{Shaded}

\end{tcolorbox}

\begin{tcolorbox}[enhanced jigsaw, colframe=quarto-callout-tip-color-frame, breakable, left=2mm, leftrule=.75mm, toptitle=1mm, rightrule=.15mm, opacitybacktitle=0.6, arc=.35mm, colback=white, title=\textcolor{quarto-callout-tip-color}{\faLightbulb}\hspace{0.5em}{Activity 2}, coltitle=black, colbacktitle=quarto-callout-tip-color!10!white, toprule=.15mm, bottomtitle=1mm, opacityback=0, titlerule=0mm, bottomrule=.15mm]

Replace \texttt{scale\_x\_continuous(limits\ =\ c(1700,\ 1815))} in the
right-hand panel above with
\texttt{coord\_cartesian(xlim\ =\ c(1700,\ 1815))}. Confirm that the
smoother shape is the same as the left panel.

\end{tcolorbox}

\bookmarksetup{startatroot}

\chapter{Grouped plots and overplotting}\label{sec-grouped-plots}

\section{Intended learning outcomes}\label{intended-learning-outcomes-4}

By the end of this chapter you will be able to:

\begin{itemize}
\tightlist
\item
  Map a categorical variable to \texttt{colour} or \texttt{fill} to
  display three variables on one plot
\item
  Distinguish setting an aesthetic \emph{inside} \texttt{aes()} from
  setting it as a fixed value \emph{outside} \texttt{aes()}
\item
  Split a plot into panels with \texttt{facet\_wrap()} and
  \texttt{facet\_grid()}
\item
  Reduce overplotting with \texttt{alpha} and \texttt{geom\_jitter()}
\end{itemize}

\section{Functions used}\label{functions-used-3}

\begin{itemize}
\tightlist
\item
  ggplot2: \texttt{aes()}, \texttt{geom\_point()},
  \texttt{geom\_jitter()}, \texttt{geom\_smooth()},
  \texttt{facet\_wrap()}, \texttt{facet\_grid()},
  \texttt{coord\_cartesian()}
\item
  \texttt{colour} and \texttt{fill} aesthetics
\end{itemize}

\section{Adding a third variable with
colour}\label{adding-a-third-variable-with-colour}

So far our plots have used at most two variables. To show how the number
who died on the Middle Passage varied across the period of British
involvement, we plot \texttt{arrival\_year} against \texttt{died} and
map \texttt{start\_port} to \texttt{colour}.

\begin{Shaded}
\begin{Highlighting}[]
\FunctionTok{ggplot}\NormalTok{(dat\_filter, }\FunctionTok{aes}\NormalTok{(arrival\_year, died, }\AttributeTok{colour =}\NormalTok{ start\_port)) }\SpecialCharTok{+}
  \FunctionTok{geom\_point}\NormalTok{(}\AttributeTok{alpha =} \FloatTok{0.4}\NormalTok{) }\SpecialCharTok{+}
  \FunctionTok{geom\_smooth}\NormalTok{() }\SpecialCharTok{+}
  \FunctionTok{scale\_colour\_viridis\_d}\NormalTok{() }\SpecialCharTok{+}
  \FunctionTok{labs}\NormalTok{(}\AttributeTok{x =} \StringTok{"Year of arrival"}\NormalTok{,}
       \AttributeTok{y =} \StringTok{"Number of deaths per voyage"}\NormalTok{,}
       \AttributeTok{colour =} \StringTok{"Port of origin"}\NormalTok{) }\SpecialCharTok{+}
  \FunctionTok{coord\_cartesian}\NormalTok{(}\AttributeTok{xlim =} \FunctionTok{c}\NormalTok{(}\DecValTok{1641}\NormalTok{, }\DecValTok{1815}\NormalTok{))}
\end{Highlighting}
\end{Shaded}

\begin{figure}[H]

\centering{

\includegraphics[width=1\linewidth,height=\textheight,keepaspectratio]{04-grouped-plots_files/figure-pdf/fig-grouped-scatter-1.pdf}

}

\caption{\label{fig-grouped-scatter}Number of deaths per voyage by year
of arrival, coloured by port of origin.}

\end{figure}%

\begin{tcolorbox}[enhanced jigsaw, colframe=quarto-callout-tip-color-frame, breakable, left=2mm, leftrule=.75mm, toptitle=1mm, rightrule=.15mm, opacitybacktitle=0.6, arc=.35mm, colback=white, title=\textcolor{quarto-callout-tip-color}{\faLightbulb}\hspace{0.5em}{Inside or outside \texttt{aes()}?}, coltitle=black, colbacktitle=quarto-callout-tip-color!10!white, toprule=.15mm, bottomtitle=1mm, opacityback=0, titlerule=0mm, bottomrule=.15mm]

\begin{itemize}
\tightlist
\item
  Inside \texttt{aes()}: the aesthetic varies with the data.
  \texttt{aes(colour\ =\ start\_port)} means each port gets a different
  colour.
\item
  Outside \texttt{aes()}: the aesthetic is a fixed value.
  \texttt{geom\_point(colour\ =\ "dodgerblue")} colours all points the
  same regardless of port.
\end{itemize}

The same rule applies to \texttt{fill}, \texttt{shape}, \texttt{size}
and \texttt{linetype}.

\end{tcolorbox}

\section{Facetting}\label{facetting}

\texttt{facet\_wrap()} splits the plot into one panel per level of a
variable. It is an alternative way to display the same three-variable
relationship and is often easier to read than overlaid colours.

\begin{Shaded}
\begin{Highlighting}[]
\FunctionTok{ggplot}\NormalTok{(dat\_filter, }\FunctionTok{aes}\NormalTok{(arrival\_year, died)) }\SpecialCharTok{+}
  \FunctionTok{geom\_point}\NormalTok{(}\FunctionTok{aes}\NormalTok{(}\AttributeTok{colour =}\NormalTok{ start\_port), }\AttributeTok{show.legend =} \ConstantTok{FALSE}\NormalTok{, }\AttributeTok{alpha =} \FloatTok{0.4}\NormalTok{) }\SpecialCharTok{+}
  \FunctionTok{geom\_smooth}\NormalTok{() }\SpecialCharTok{+}
  \FunctionTok{scale\_colour\_viridis\_d}\NormalTok{() }\SpecialCharTok{+}
  \FunctionTok{facet\_wrap}\NormalTok{(}\SpecialCharTok{\textasciitilde{}}\NormalTok{ start\_port, }\AttributeTok{nrow =} \DecValTok{3}\NormalTok{) }\SpecialCharTok{+}
  \FunctionTok{labs}\NormalTok{(}\AttributeTok{x =} \StringTok{"Year of arrival"}\NormalTok{,}
       \AttributeTok{y =} \StringTok{"Number of deaths per voyage"}\NormalTok{,}
       \AttributeTok{title =} \StringTok{"Deaths during the Middle Passage, by port of origin"}\NormalTok{) }\SpecialCharTok{+}
  \FunctionTok{coord\_cartesian}\NormalTok{(}\AttributeTok{xlim =} \FunctionTok{c}\NormalTok{(}\DecValTok{1641}\NormalTok{, }\DecValTok{1815}\NormalTok{))}
\end{Highlighting}
\end{Shaded}

\begin{figure}[H]

\centering{

\includegraphics[width=1\linewidth,height=\textheight,keepaspectratio]{04-grouped-plots_files/figure-pdf/fig-facet-wrap-1.pdf}

}

\caption{\label{fig-facet-wrap}Same data, faceted by port of origin.}

\end{figure}%

The base mapping for the smoother uses no \texttt{colour} mapping, so a
single smoother is drawn per facet. Inside \texttt{geom\_point()} we map
\texttt{colour\ =\ start\_port} so the points within each facet are
coloured by port even though there is only one port per panel; this is
purely decorative.

\begin{itemize}
\tightlist
\item
  Try changing \texttt{nrow\ =\ 3} to \texttt{nrow\ =\ 1} and rendering.
  What happens?
\end{itemize}

\begin{tcolorbox}[enhanced jigsaw, colframe=quarto-callout-tip-color-frame, breakable, left=2mm, leftrule=.75mm, toptitle=1mm, rightrule=.15mm, opacitybacktitle=0.6, arc=.35mm, colback=white, title=\textcolor{quarto-callout-tip-color}{\faLightbulb}\hspace{0.5em}{Activity}, coltitle=black, colbacktitle=quarto-callout-tip-color!10!white, toprule=.15mm, bottomtitle=1mm, opacityback=0, titlerule=0mm, bottomrule=.15mm]

Replace
\texttt{facet\_wrap(\textasciitilde{}\ start\_port,\ nrow\ =\ 3)} with
\texttt{facet\_grid(start\_port\ \textasciitilde{}\ .)}. Read the help
with \texttt{?facet\_grid} to see when one is preferred over the other.

\end{tcolorbox}

\section{Overplotting}\label{sec-overplotting}

The earliest scatterplots in Chapter~\ref{sec-basic-plots} had so many
overlapping points that the density of the data was hidden. Two
techniques help.

\texttt{alpha} makes each point semi-transparent, so dense regions
appear darker.

\texttt{geom\_jitter()} adds a small random offset, so points stacked on
the same value are visible separately.

\begin{Shaded}
\begin{Highlighting}[]
\NormalTok{opaque }\OtherTok{\textless{}{-}} \FunctionTok{ggplot}\NormalTok{(dat\_filter, }\FunctionTok{aes}\NormalTok{(start\_port, mortality)) }\SpecialCharTok{+}
  \FunctionTok{geom\_point}\NormalTok{() }\SpecialCharTok{+}
  \FunctionTok{labs}\NormalTok{(}\AttributeTok{title =} \StringTok{"alpha = 1 (default)"}\NormalTok{)}

\NormalTok{transparent }\OtherTok{\textless{}{-}} \FunctionTok{ggplot}\NormalTok{(dat\_filter, }\FunctionTok{aes}\NormalTok{(start\_port, mortality)) }\SpecialCharTok{+}
  \FunctionTok{geom\_point}\NormalTok{(}\AttributeTok{alpha =} \FloatTok{0.1}\NormalTok{) }\SpecialCharTok{+}
  \FunctionTok{labs}\NormalTok{(}\AttributeTok{title =} \StringTok{"alpha = 0.1"}\NormalTok{)}

\NormalTok{jittered }\OtherTok{\textless{}{-}} \FunctionTok{ggplot}\NormalTok{(dat\_filter, }\FunctionTok{aes}\NormalTok{(start\_port, mortality)) }\SpecialCharTok{+}
  \FunctionTok{geom\_jitter}\NormalTok{(}\AttributeTok{alpha =} \FloatTok{0.2}\NormalTok{, }\AttributeTok{width =} \FloatTok{0.3}\NormalTok{, }\AttributeTok{height =} \DecValTok{0}\NormalTok{) }\SpecialCharTok{+}
  \FunctionTok{labs}\NormalTok{(}\AttributeTok{title =} \StringTok{"geom\_jitter()"}\NormalTok{)}

\NormalTok{opaque }\SpecialCharTok{+}\NormalTok{ transparent }\SpecialCharTok{+}\NormalTok{ jittered}
\end{Highlighting}
\end{Shaded}

\begin{figure}[H]

\centering{

\includegraphics[width=1\linewidth,height=\textheight,keepaspectratio]{04-grouped-plots_files/figure-pdf/fig-overplotting-1.pdf}

}

\caption{\label{fig-overplotting}Three approaches to overplotting.}

\end{figure}%

\begin{tcolorbox}[enhanced jigsaw, colframe=quarto-callout-warning-color-frame, breakable, left=2mm, leftrule=.75mm, toptitle=1mm, rightrule=.15mm, opacitybacktitle=0.6, arc=.35mm, colback=white, title=\textcolor{quarto-callout-warning-color}{\faExclamationTriangle}\hspace{0.5em}{Warning}, coltitle=black, colbacktitle=quarto-callout-warning-color!10!white, toprule=.15mm, bottomtitle=1mm, opacityback=0, titlerule=0mm, bottomrule=.15mm]

\texttt{geom\_jitter()} moves points slightly. That is fine for showing
density on a categorical axis, but if you jitter on a continuous axis
you are misrepresenting the data. Set the \texttt{width} and
\texttt{height} arguments deliberately. Here we use
\texttt{height\ =\ 0} so vertical positions are unchanged.

\end{tcolorbox}

\bookmarksetup{startatroot}

\chapter{Reshaping data and combining plots}\label{sec-reshaping}

\section{Intended learning outcomes}\label{intended-learning-outcomes-5}

By the end of this chapter you will be able to:

\begin{itemize}
\tightlist
\item
  Reshape data from wide to long with \texttt{pivot\_longer()}
\item
  Plot grouped distributions using one mapping per categorical column
  rather than one column per category
\item
  Order categorical axes deliberately with \texttt{factor()} or
  \texttt{fct\_relevel()}
\item
  Combine multiple plots into a single figure with the
  \texttt{patchwork} package
\end{itemize}

\section{Functions used}\label{functions-used-4}

\begin{itemize}
\tightlist
\item
  tidyr: \texttt{pivot\_longer()}, \texttt{drop\_na()}
\item
  forcats: \texttt{fct\_relevel()}
\item
  ggplot2: \texttt{geom\_boxplot()}
\item
  patchwork: \texttt{+}, \texttt{/}, \texttt{plot\_layout()},
  \texttt{plot\_annotation()}
\end{itemize}

\section{From wide to long}\label{sec-wide-to-long}

The dataset has three columns recording the demographic make-up of each
voyage. \texttt{men}, \texttt{women} and \texttt{children} are each a
proportion of the people embarked. Most voyages in the database do not
record this breakdown, so these three columns are missing for the great
majority of rows. Before reshaping, we drop the voyages with no
demographic data using \texttt{drop\_na()}, which removes any row that
is missing a value in one of the named columns.

\begin{Shaded}
\begin{Highlighting}[]
\NormalTok{dat\_demographic }\OtherTok{\textless{}{-}}\NormalTok{ dat\_filter }\SpecialCharTok{|\textgreater{}}
  \FunctionTok{drop\_na}\NormalTok{(men, women, children)}

\FunctionTok{head}\NormalTok{(dat\_demographic[, }\FunctionTok{c}\NormalTok{(}\StringTok{"start\_port"}\NormalTok{, }\StringTok{"men"}\NormalTok{, }\StringTok{"women"}\NormalTok{, }\StringTok{"children"}\NormalTok{)])}
\end{Highlighting}
\end{Shaded}

\begin{verbatim}
# A tibble: 6 x 4
  start_port   men women children
  <fct>      <dbl> <dbl>    <dbl>
1 Bristol    0.384 0.232    0.384
2 Bristol    0.392 0.256    0.352
3 Bristol    0.488 0.175    0.337
4 Bristol    0.269 0.223    0.508
5 Bristol    0.149 0.131    0.720
6 Bristol    0.477 0.185    0.337
\end{verbatim}

This is the \emph{wide} form. Each demographic category has its own
column. To plot them together we need the \emph{long} form, where one
column holds the category and another holds the value.

\texttt{pivot\_longer()} does the conversion. We tell it which columns
to gather (\texttt{cols}), the name of the new column that will hold the
original column names (\texttt{names\_to}) and the name of the new
column that will hold the values (\texttt{values\_to}).

\begin{Shaded}
\begin{Highlighting}[]
\NormalTok{dat\_long }\OtherTok{\textless{}{-}}\NormalTok{ dat\_demographic }\SpecialCharTok{|\textgreater{}}
  \FunctionTok{pivot\_longer}\NormalTok{(}\AttributeTok{cols =} \FunctionTok{c}\NormalTok{(men, women, children),}
               \AttributeTok{names\_to =} \StringTok{"demographic"}\NormalTok{,}
               \AttributeTok{values\_to =} \StringTok{"proportion"}\NormalTok{)}

\FunctionTok{head}\NormalTok{(dat\_long[, }\FunctionTok{c}\NormalTok{(}\StringTok{"start\_port"}\NormalTok{, }\StringTok{"demographic"}\NormalTok{, }\StringTok{"proportion"}\NormalTok{)])}
\end{Highlighting}
\end{Shaded}

\begin{verbatim}
# A tibble: 6 x 3
  start_port demographic proportion
  <fct>      <chr>            <dbl>
1 Bristol    men              0.384
2 Bristol    women            0.232
3 Bristol    children         0.384
4 Bristol    men              0.392
5 Bristol    women            0.256
6 Bristol    children         0.352
\end{verbatim}

Each voyage now appears three times, once per demographic category. The
column is called \texttt{demographic} because it holds the categories
\emph{man}, \emph{woman} and \emph{child}. The values are proportions of
the people embarked.

\begin{tcolorbox}[enhanced jigsaw, colframe=quarto-callout-note-color-frame, breakable, left=2mm, leftrule=.75mm, toptitle=1mm, rightrule=.15mm, opacitybacktitle=0.6, arc=.35mm, colback=white, title=\textcolor{quarto-callout-note-color}{\faInfo}\hspace{0.5em}{Note}, coltitle=black, colbacktitle=quarto-callout-note-color!10!white, toprule=.15mm, bottomtitle=1mm, opacityback=0, titlerule=0mm, bottomrule=.15mm]

The demographic figures describe only the voyages that recorded them, a
minority of the total. The original column names \texttt{men},
\texttt{women} and \texttt{children} are inherited from the source
database. In the reshaped data, the column \texttt{demographic} holds
those labels as values.

\end{tcolorbox}

\section{Ordering categorical axes}\label{ordering-categorical-axes}

By default, ggplot will plot \texttt{demographic} in alphabetical order,
which puts \texttt{children} first. We probably want \texttt{men},
\texttt{women}, \texttt{children} or the other way around.
\texttt{fct\_relevel()} from \texttt{forcats} is the easiest way.

\begin{Shaded}
\begin{Highlighting}[]
\NormalTok{dat\_long }\OtherTok{\textless{}{-}}\NormalTok{ dat\_long }\SpecialCharTok{|\textgreater{}}
  \FunctionTok{mutate}\NormalTok{(}\AttributeTok{demographic =} \FunctionTok{fct\_relevel}\NormalTok{(demographic, }\StringTok{"men"}\NormalTok{, }\StringTok{"women"}\NormalTok{, }\StringTok{"children"}\NormalTok{))}
\end{Highlighting}
\end{Shaded}

\section{Grouped boxplots}\label{grouped-boxplots}

We can now show the distribution of the demographic proportions, broken
down by port of origin, all on one plot.

\begin{Shaded}
\begin{Highlighting}[]
\FunctionTok{ggplot}\NormalTok{(dat\_long, }\FunctionTok{aes}\NormalTok{(}\AttributeTok{x =}\NormalTok{ demographic, }\AttributeTok{y =}\NormalTok{ proportion, }\AttributeTok{fill =}\NormalTok{ start\_port)) }\SpecialCharTok{+}
  \FunctionTok{geom\_boxplot}\NormalTok{(}\AttributeTok{alpha =} \FloatTok{0.6}\NormalTok{) }\SpecialCharTok{+}
  \FunctionTok{scale\_fill\_viridis\_d}\NormalTok{(}\AttributeTok{option =} \StringTok{"E"}\NormalTok{) }\SpecialCharTok{+}
  \FunctionTok{labs}\NormalTok{(}\AttributeTok{x =} \StringTok{"Demographic category"}\NormalTok{,}
       \AttributeTok{y =} \StringTok{"Proportion of those embarked"}\NormalTok{,}
       \AttributeTok{fill =} \StringTok{"Port of origin"}\NormalTok{)}
\end{Highlighting}
\end{Shaded}

\begin{figure}[H]

\centering{

\includegraphics[width=1\linewidth,height=\textheight,keepaspectratio]{05-reshaping_files/figure-pdf/fig-grouped-box-1.pdf}

}

\caption{\label{fig-grouped-box}Demographic composition of voyages from
each port.}

\end{figure}%

Men make up the majority of those embarked on every type of voyage, but
the proportion of women and children varies markedly by port.

\begin{tcolorbox}[enhanced jigsaw, colframe=quarto-callout-tip-color-frame, breakable, left=2mm, leftrule=.75mm, toptitle=1mm, rightrule=.15mm, opacitybacktitle=0.6, arc=.35mm, colback=white, title=\textcolor{quarto-callout-tip-color}{\faLightbulb}\hspace{0.5em}{Activity}, coltitle=black, colbacktitle=quarto-callout-tip-color!10!white, toprule=.15mm, bottomtitle=1mm, opacityback=0, titlerule=0mm, bottomrule=.15mm]

Change the value of \texttt{alpha} and observe what happens. Try
\texttt{alpha\ =\ 0.2} and \texttt{alpha\ =\ 1}.

\end{tcolorbox}

\section{Piping straight into ggplot}\label{piping-straight-into-ggplot}

You can combine \texttt{filter()} and \texttt{ggplot()} in a single pipe
rather than creating an intermediate object. The example below takes the
long-format data, filters to voyages that disembarked in the Americas,
and plots the result.

\begin{Shaded}
\begin{Highlighting}[]
\NormalTok{dat\_long }\SpecialCharTok{|\textgreater{}}
  \FunctionTok{filter}\NormalTok{(outcome }\SpecialCharTok{==} \StringTok{"Slaves disembarked in Americas"}\NormalTok{) }\SpecialCharTok{|\textgreater{}}
  \FunctionTok{ggplot}\NormalTok{(}\FunctionTok{aes}\NormalTok{(}\AttributeTok{x =}\NormalTok{ demographic, }\AttributeTok{y =}\NormalTok{ proportion, }\AttributeTok{fill =}\NormalTok{ demographic)) }\SpecialCharTok{+}
  \FunctionTok{geom\_boxplot}\NormalTok{(}\AttributeTok{show.legend =} \ConstantTok{FALSE}\NormalTok{) }\SpecialCharTok{+}
  \FunctionTok{scale\_fill\_viridis\_d}\NormalTok{(}\AttributeTok{option =} \StringTok{"E"}\NormalTok{) }\SpecialCharTok{+}
  \FunctionTok{labs}\NormalTok{(}\AttributeTok{x =} \ConstantTok{NULL}\NormalTok{,}
       \AttributeTok{y =} \StringTok{"Proportion of those embarked"}\NormalTok{,}
       \AttributeTok{title =} \StringTok{"Demographic composition of voyages disembarking in the Americas"}\NormalTok{)}
\end{Highlighting}
\end{Shaded}

\begin{figure}[H]

\centering{

\includegraphics[width=1\linewidth,height=\textheight,keepaspectratio]{05-reshaping_files/figure-pdf/fig-filtered-box-1.pdf}

}

\caption{\label{fig-filtered-box}Demographic composition of voyages that
disembarked in the Americas.}

\end{figure}%

\section{Combining plots with patchwork}\label{sec-patchwork}

\texttt{patchwork} lets you arrange \texttt{ggplot} objects into
multi-panel figures. The syntax is intuitive: \texttt{+} puts plots side
by side, \texttt{/} stacks them vertically, and \texttt{plot\_layout()}
controls the arrangement.

First, save two plots as objects.

\begin{Shaded}
\begin{Highlighting}[]
\NormalTok{men\_plot }\OtherTok{\textless{}{-}}\NormalTok{ dat\_long }\SpecialCharTok{|\textgreater{}}
  \FunctionTok{filter}\NormalTok{(demographic }\SpecialCharTok{==} \StringTok{"men"}\NormalTok{) }\SpecialCharTok{|\textgreater{}}
  \FunctionTok{ggplot}\NormalTok{(}\FunctionTok{aes}\NormalTok{(start\_port, proportion, }\AttributeTok{fill =}\NormalTok{ start\_port)) }\SpecialCharTok{+}
  \FunctionTok{geom\_boxplot}\NormalTok{(}\AttributeTok{show.legend =} \ConstantTok{FALSE}\NormalTok{, }\AttributeTok{alpha =} \FloatTok{0.6}\NormalTok{) }\SpecialCharTok{+}
  \FunctionTok{scale\_fill\_viridis\_d}\NormalTok{(}\AttributeTok{option =} \StringTok{"E"}\NormalTok{) }\SpecialCharTok{+}
  \FunctionTok{labs}\NormalTok{(}\AttributeTok{x =} \ConstantTok{NULL}\NormalTok{, }\AttributeTok{y =} \StringTok{"Proportion men"}\NormalTok{)}

\NormalTok{children\_plot }\OtherTok{\textless{}{-}}\NormalTok{ dat\_long }\SpecialCharTok{|\textgreater{}}
  \FunctionTok{filter}\NormalTok{(demographic }\SpecialCharTok{==} \StringTok{"children"}\NormalTok{) }\SpecialCharTok{|\textgreater{}}
  \FunctionTok{ggplot}\NormalTok{(}\FunctionTok{aes}\NormalTok{(start\_port, proportion, }\AttributeTok{fill =}\NormalTok{ start\_port)) }\SpecialCharTok{+}
  \FunctionTok{geom\_boxplot}\NormalTok{(}\AttributeTok{show.legend =} \ConstantTok{FALSE}\NormalTok{, }\AttributeTok{alpha =} \FloatTok{0.6}\NormalTok{) }\SpecialCharTok{+}
  \FunctionTok{scale\_fill\_viridis\_d}\NormalTok{(}\AttributeTok{option =} \StringTok{"E"}\NormalTok{) }\SpecialCharTok{+}
  \FunctionTok{labs}\NormalTok{(}\AttributeTok{x =} \ConstantTok{NULL}\NormalTok{, }\AttributeTok{y =} \StringTok{"Proportion children"}\NormalTok{)}
\end{Highlighting}
\end{Shaded}

Then combine them.

\begin{Shaded}
\begin{Highlighting}[]
\NormalTok{men\_plot }\SpecialCharTok{+}\NormalTok{ children\_plot}
\end{Highlighting}
\end{Shaded}

\begin{figure}[H]

\centering{

\includegraphics[width=1\linewidth,height=\textheight,keepaspectratio]{05-reshaping_files/figure-pdf/fig-patchwork-side-1.pdf}

}

\caption{\label{fig-patchwork-side}Two boxplots side by side with
patchwork.}

\end{figure}%

Stacked vertically.

\begin{Shaded}
\begin{Highlighting}[]
\NormalTok{men\_plot }\SpecialCharTok{/}\NormalTok{ children\_plot}
\end{Highlighting}
\end{Shaded}

\begin{figure}[H]

\centering{

\includegraphics[width=1\linewidth,height=\textheight,keepaspectratio]{05-reshaping_files/figure-pdf/fig-patchwork-stack-1.pdf}

}

\caption{\label{fig-patchwork-stack}The same two boxplots stacked
vertically.}

\end{figure}%

\texttt{plot\_annotation()} adds a single overall title across multiple
panels.

\begin{Shaded}
\begin{Highlighting}[]
\NormalTok{(men\_plot }\SpecialCharTok{+}\NormalTok{ children\_plot) }\SpecialCharTok{+}
  \FunctionTok{plot\_annotation}\NormalTok{(}\AttributeTok{title =} \StringTok{"Demographic composition by port of origin"}\NormalTok{)}
\end{Highlighting}
\end{Shaded}

\begin{figure}[H]

\centering{

\includegraphics[width=1\linewidth,height=\textheight,keepaspectratio]{05-reshaping_files/figure-pdf/fig-patchwork-annotated-1.pdf}

}

\caption{\label{fig-patchwork-annotated}Combined figure with a shared
title.}

\end{figure}%

See the \href{https://patchwork.data-imaginist.com/}{\texttt{patchwork}
documentation} for more elaborate layouts, shared legends and
\texttt{tag\_levels} for automatic panel labels (A, B, C).

\section{Activities}\label{activities-3}

\begin{tcolorbox}[enhanced jigsaw, colframe=quarto-callout-tip-color-frame, breakable, left=2mm, leftrule=.75mm, toptitle=1mm, rightrule=.15mm, opacitybacktitle=0.6, arc=.35mm, colback=white, title=\textcolor{quarto-callout-tip-color}{\faLightbulb}\hspace{0.5em}{Activity 1}, coltitle=black, colbacktitle=quarto-callout-tip-color!10!white, toprule=.15mm, bottomtitle=1mm, opacityback=0, titlerule=0mm, bottomrule=.15mm]

Build a third plot showing the proportion of women by port, and combine
all three plots (\texttt{men\_plot}, \texttt{women\_plot},
\texttt{children\_plot}) in a single row.

\end{tcolorbox}

\begin{tcolorbox}[enhanced jigsaw, colframe=quarto-callout-note-color-frame, breakable, left=2mm, leftrule=.75mm, toptitle=1mm, rightrule=.15mm, opacitybacktitle=0.6, arc=.35mm, colback=white, title=\textcolor{quarto-callout-note-color}{\faInfo}\hspace{0.5em}{Solution}, coltitle=black, colbacktitle=quarto-callout-note-color!10!white, toprule=.15mm, bottomtitle=1mm, opacityback=0, titlerule=0mm, bottomrule=.15mm]

\begin{Shaded}
\begin{Highlighting}[]
\NormalTok{women\_plot }\OtherTok{\textless{}{-}}\NormalTok{ dat\_long }\SpecialCharTok{|\textgreater{}}
  \FunctionTok{filter}\NormalTok{(demographic }\SpecialCharTok{==} \StringTok{"women"}\NormalTok{) }\SpecialCharTok{|\textgreater{}}
  \FunctionTok{ggplot}\NormalTok{(}\FunctionTok{aes}\NormalTok{(start\_port, proportion, }\AttributeTok{fill =}\NormalTok{ start\_port)) }\SpecialCharTok{+}
  \FunctionTok{geom\_boxplot}\NormalTok{(}\AttributeTok{show.legend =} \ConstantTok{FALSE}\NormalTok{, }\AttributeTok{alpha =} \FloatTok{0.6}\NormalTok{) }\SpecialCharTok{+}
  \FunctionTok{scale\_fill\_viridis\_d}\NormalTok{(}\AttributeTok{option =} \StringTok{"E"}\NormalTok{) }\SpecialCharTok{+}
  \FunctionTok{labs}\NormalTok{(}\AttributeTok{x =} \ConstantTok{NULL}\NormalTok{, }\AttributeTok{y =} \StringTok{"Proportion women"}\NormalTok{)}

\NormalTok{men\_plot }\SpecialCharTok{+}\NormalTok{ women\_plot }\SpecialCharTok{+}\NormalTok{ children\_plot}
\end{Highlighting}
\end{Shaded}

\end{tcolorbox}

\begin{tcolorbox}[enhanced jigsaw, colframe=quarto-callout-tip-color-frame, breakable, left=2mm, leftrule=.75mm, toptitle=1mm, rightrule=.15mm, opacitybacktitle=0.6, arc=.35mm, colback=white, title=\textcolor{quarto-callout-tip-color}{\faLightbulb}\hspace{0.5em}{Activity 2}, coltitle=black, colbacktitle=quarto-callout-tip-color!10!white, toprule=.15mm, bottomtitle=1mm, opacityback=0, titlerule=0mm, bottomrule=.15mm]

Reshape the data back into wide format using \texttt{pivot\_wider()}.
(Hint: it takes \texttt{names\_from} and \texttt{values\_from}
arguments.)

\end{tcolorbox}

\bookmarksetup{startatroot}

\chapter{Saving and sharing plots}\label{sec-saving}

\section{Intended learning outcomes}\label{intended-learning-outcomes-6}

By the end of this chapter you will be able to:

\begin{itemize}
\tightlist
\item
  Save a plot to a file with \texttt{ggsave()}
\item
  Choose appropriate dimensions and resolution for print or web
\item
  Write alt text for screen-reader accessibility
\item
  Find further resources to extend your skills
\end{itemize}

\section{Functions used}\label{functions-used-5}

\begin{itemize}
\tightlist
\item
  ggplot2: \texttt{ggsave()}
\item
  Chunk options: \texttt{\#\textbar{}\ fig-cap:},
  \texttt{\#\textbar{}\ fig-alt:}, \texttt{\#\textbar{}\ fig-width:},
  \texttt{\#\textbar{}\ fig-height:}
\end{itemize}

\section{\texorpdfstring{\texttt{ggsave()}
basics}{ggsave() basics}}\label{ggsave-basics}

\texttt{ggsave()} saves the most recently produced plot to a file, with
the format inferred from the file extension. Supported formats include
\texttt{.png}, \texttt{.jpeg}, \texttt{.pdf}, \texttt{.svg},
\texttt{.tiff}, \texttt{.bmp}, \texttt{.eps} and \texttt{.wmf}.

\begin{Shaded}
\begin{Highlighting}[]
\FunctionTok{ggsave}\NormalTok{(}\StringTok{"last\_plot.png"}\NormalTok{)}
\end{Highlighting}
\end{Shaded}

To save a specific plot rather than the last one drawn, pass it as the
\texttt{plot} argument.

\begin{Shaded}
\begin{Highlighting}[]
\NormalTok{my\_plot }\OtherTok{\textless{}{-}} \FunctionTok{ggplot}\NormalTok{(dat\_filter, }\FunctionTok{aes}\NormalTok{(}\AttributeTok{x =}\NormalTok{ start\_port, }\AttributeTok{fill =}\NormalTok{ start\_port)) }\SpecialCharTok{+}
  \FunctionTok{geom\_bar}\NormalTok{(}\AttributeTok{show.legend =} \ConstantTok{FALSE}\NormalTok{)}

\FunctionTok{ggsave}\NormalTok{(}\StringTok{"my\_plot.png"}\NormalTok{, }\AttributeTok{plot =}\NormalTok{ my\_plot)}
\end{Highlighting}
\end{Shaded}

\section{Controlling size and
resolution}\label{controlling-size-and-resolution}

By default \texttt{ggsave()} uses the dimensions of the active graphics
device, which usually means whatever your RStudio plot pane happens to
be. For consistent output, set \texttt{width}, \texttt{height} and
\texttt{units} explicitly.

\begin{Shaded}
\begin{Highlighting}[]
\FunctionTok{ggsave}\NormalTok{(}\StringTok{"my\_plot.png"}\NormalTok{,}
       \AttributeTok{plot =}\NormalTok{ my\_plot,}
       \AttributeTok{width  =} \DecValTok{16}\NormalTok{,}
       \AttributeTok{height =} \DecValTok{10}\NormalTok{,}
       \AttributeTok{units  =} \StringTok{"cm"}\NormalTok{,}
       \AttributeTok{dpi    =} \DecValTok{300}\NormalTok{)}
\end{Highlighting}
\end{Shaded}

\begin{tcolorbox}[enhanced jigsaw, colframe=quarto-callout-tip-color-frame, breakable, left=2mm, leftrule=.75mm, toptitle=1mm, rightrule=.15mm, opacitybacktitle=0.6, arc=.35mm, colback=white, title=\textcolor{quarto-callout-tip-color}{\faLightbulb}\hspace{0.5em}{Tip}, coltitle=black, colbacktitle=quarto-callout-tip-color!10!white, toprule=.15mm, bottomtitle=1mm, opacityback=0, titlerule=0mm, bottomrule=.15mm]

\begin{itemize}
\tightlist
\item
  For posters and prints, use centimetres or inches and a high
  \texttt{dpi} (300+).
\item
  For web display, use pixels (\texttt{units\ =\ "px"}) and
  \texttt{dpi\ =\ 72} to \texttt{96}.
\item
  For papers, journals typically want a vector format such as
  \texttt{.pdf} or \texttt{.svg} so the figure scales without
  pixelation. Check the journal's submission guide.
\end{itemize}

\end{tcolorbox}

\section{Accessibility: alt text}\label{accessibility-alt-text}

A plot that is meaningful to a sighted reader is invisible to a screen
reader unless you add alt text. In Quarto, \texttt{fig-alt} is a chunk
option that becomes the \texttt{alt} attribute on the rendered image.

\begin{Shaded}
\begin{Highlighting}[]
\InformationTok{\textasciigrave{}\textasciigrave{}\textasciigrave{}\{r\}}
\CommentTok{\#| label: fig{-}mortality}
\CommentTok{\#| fig{-}cap: "Mortality rate by port of origin."}
\CommentTok{\#| fig{-}alt: "Boxplots of mortality rate per voyage for Liverpool, London and Bristol. All three medians lie between 0.1 and 0.2, with outliers approaching 1.0."}

\FunctionTok{ggplot}\NormalTok{(dat\_filter, }\FunctionTok{aes}\NormalTok{(start\_port, mortality)) }\SpecialCharTok{+}
  \FunctionTok{geom\_boxplot}\NormalTok{()}
\InformationTok{\textasciigrave{}\textasciigrave{}\textasciigrave{}}
\end{Highlighting}
\end{Shaded}

A useful alt text describes the data first (what is on each axis, what
the broad pattern is) and the decoration second (colours, themes). It is
for someone who cannot see the image, not someone who is choosing
whether to look at it.

\begin{tcolorbox}[enhanced jigsaw, colframe=quarto-callout-warning-color-frame, breakable, left=2mm, leftrule=.75mm, toptitle=1mm, rightrule=.15mm, opacitybacktitle=0.6, arc=.35mm, colback=white, title=\textcolor{quarto-callout-warning-color}{\faExclamationTriangle}\hspace{0.5em}{Warning}, coltitle=black, colbacktitle=quarto-callout-warning-color!10!white, toprule=.15mm, bottomtitle=1mm, opacityback=0, titlerule=0mm, bottomrule=.15mm]

Setting \texttt{caption\ =} inside \texttt{labs()} puts text inside the
rendered plot, where a screen reader cannot reach it. Use the
\texttt{fig-cap} chunk option instead, which lives in the surrounding
HTML and is read aloud.

\end{tcolorbox}

\section{Where to go next}\label{where-to-go-next}

\begin{itemize}
\tightlist
\item
  Chapter~\ref{sec-getting-started} covered the basics. For a deeper
  treatment of the same material, see
  \href{https://psyteachr.github.io/introdataviz/}{Data Visualisation
  Using R, for Researchers Who Don't Use R} (Nordmann et al., 2022).
\item
  For a broader course covering Quarto, data wrangling, joins and
  reproducible reports, see
  \href{https://psyteachr.github.io/ads-v3/}{Applied Data Skills}
  (Nordmann \& DeBruine, 2025).
\item
  For a reference card of \texttt{ggplot2} functions, the
  \href{https://rstudio.github.io/cheatsheets/html/data-visualization.html}{official
  cheat sheet} is excellent.
\item
  The \href{https://r-graph-gallery.com/}{R Graph Gallery} is a
  searchable collection of every chart type imaginable, with the code to
  produce each one.
\item
  For statistical concepts and worked examples, see
  \href{https://r4ds.hadley.nz/}{R for Data Science} (Wickham et al.,
  2023), which is freely available online.
\end{itemize}

\section{A closing note}\label{a-closing-note}

This tutorial is built around a dataset documenting an atrocity. The
plots you have produced quantify the scale and mortality of the
trans-Atlantic trade in enslaved people, and the prominence of British
ports in it. Plots are not neutral, however carefully you draw them.
Choices about axes, colours, what to aggregate and what to leave
granular all shape the story a figure tells. Bringing those choices into
conscious view is part of what makes data visualisation a worthwhile
skill.

\cleardoublepage
\phantomsection
\addcontentsline{toc}{part}{Appendices}
\appendix

\chapter{Installing R and RStudio}\label{sec-installing-r}

This tutorial assumes you already have R, RStudio and the
\texttt{tidyverse} installed. If you are starting from scratch, the most
up-to-date instructions are in chapter 1 of \emph{Applied Data Skills}:

\begin{itemize}
\tightlist
\item
  \href{https://psyteachr.github.io/ads-v3/01-intro.html\#sec-install-r}{Installing
  R and RStudio}
\item
  \href{https://psyteachr.github.io/ads-v3/01-intro.html\#sec-install-package}{Installing
  packages}
\item
  \href{https://psyteachr.github.io/ads-v3/01-intro.html\#sec-rproj}{Setting
  up a project}
\end{itemize}

In short:

\begin{enumerate}
\def\labelenumi{\arabic{enumi}.}
\tightlist
\item
  Install R from \url{https://cran.r-project.org/} (choose the version
  for your operating system).
\item
  Install RStudio Desktop from
  \url{https://posit.co/download/rstudio-desktop/}. RStudio is a
  development environment for R; it is free for individual use.
\item
  Install Quarto from \url{https://quarto.org/docs/get-started/}. Quarto
  is what renders this book and what you will use to produce reports of
  your own.
\item
  In RStudio, run
  \texttt{install.packages(c("tidyverse",\ "patchwork"))} in the
  console. This installs the packages used throughout this tutorial.
\end{enumerate}

You only need to install each package once per machine. You then load it
with \texttt{library()} at the top of each new session, as we do in
Chapter~\ref{sec-getting-started}.

\begin{tcolorbox}[enhanced jigsaw, colframe=quarto-callout-tip-color-frame, breakable, left=2mm, leftrule=.75mm, toptitle=1mm, rightrule=.15mm, opacitybacktitle=0.6, arc=.35mm, colback=white, title=\textcolor{quarto-callout-tip-color}{\faLightbulb}\hspace{0.5em}{Tip}, coltitle=black, colbacktitle=quarto-callout-tip-color!10!white, toprule=.15mm, bottomtitle=1mm, opacityback=0, titlerule=0mm, bottomrule=.15mm]

If you cannot install R locally, \href{https://posit.cloud/}{Posit
Cloud} gives you a free RStudio environment in the browser with no
installation required. It is sufficient for working through this
tutorial.

\end{tcolorbox}

\chapter{Glossary}\label{sec-glossary}

\begin{verbatim}
data frame with 0 columns and 0 rows
\end{verbatim}

\chapter{Further resources}\label{sec-resources}

\section{Companion tutorials}\label{companion-tutorials}

\begin{itemize}
\tightlist
\item
  \href{https://psyteachr.github.io/introdataviz/}{Data Visualisation
  Using R, for Researchers Who Don't Use R}. Nordmann, McAleer, Toivo,
  Paterson \& DeBruine (2022). The published article that accompanies
  this is at \url{https://doi.org/10.1177/25152459221074654}.
\item
  \href{https://psyteachr.github.io/ads-v3/}{Applied Data Skills}.
  Nordmann \& DeBruine (2025). A 10-week course covering Quarto reports,
  data visualisation, summaries, joins, tidying and wrangling.
\end{itemize}

\section{Reference}\label{reference}

\begin{itemize}
\tightlist
\item
  The
  \href{https://rstudio.github.io/cheatsheets/html/data-visualization.html}{ggplot2
  cheat sheet} is the single most useful one-page summary.
\item
  The \href{https://ggplot2.tidyverse.org/reference/index.html}{ggplot2
  reference site} has the canonical documentation for every function
  with worked examples.
\item
  The
  \href{https://rstudio.github.io/cheatsheets/html/data-transformation.html}{dplyr
  cheat sheet} and
  \href{https://rstudio.github.io/cheatsheets/html/data-import.html}{tidyr
  cheat sheet} cover wrangling and reshaping.
\item
  The \href{https://r-graph-gallery.com/}{R Graph Gallery} is a
  searchable collection of nearly every chart type, with code.
\end{itemize}

\section{Books}\label{books}

\begin{itemize}
\tightlist
\item
  \href{https://r4ds.hadley.nz/}{R for Data Science} (Wickham,
  Çetinkaya-Rundel \& Grolemund). Free online. The standard introduction
  to \texttt{tidyverse} for data analysis.
\item
  \href{https://ggplot2-book.org/}{ggplot2: Elegant Graphics for Data
  Analysis} (Wickham et al.). Free online. The definitive
  \texttt{ggplot2} book.
\item
  \href{https://clauswilke.com/dataviz/}{Fundamentals of Data
  Visualization} (Wilke). Free online. Principles rather than code;
  pairs well with the others.
\end{itemize}

\section{Quarto}\label{quarto}

\begin{itemize}
\tightlist
\item
  \href{https://quarto.org/}{Quarto.org} is the home of the format used
  by this book.
\item
  The \href{https://quarto.org/docs/books/}{Quarto guide for books}
  covers building multi-chapter projects like this one.
\item
  \href{https://psyteachr.github.io/ads-v3/02-reports.html}{\texttt{ads-v3}
  chapter 2} is a gentle introduction to writing Quarto reports.
\end{itemize}

\section{Data sources}\label{data-sources}

\begin{itemize}
\tightlist
\item
  \href{https://www.slavevoyages.org/}{Slave Voyages}. The
  Trans-Atlantic and Intra-American Slave Trade Databases. Includes
  documentation, maps and an estimates interface.
\end{itemize}

\chapter*{References}\label{references}
\addcontentsline{toc}{chapter}{References}

\markboth{References}{References}

\phantomsection\label{refs}
\begin{CSLReferences}{1}{0}
\bibitem[\citeproctext]{ref-nordmann2025ads}
Nordmann, E., \& DeBruine, L. M. (2025). \emph{Applied data skills}.
\url{https://psyteachr.github.io/ads-v3/}

\bibitem[\citeproctext]{ref-nordmann2022}
Nordmann, E., McAleer, P., Toivo, W., Paterson, H., \& DeBruine, L. M.
(2022). Data visualization using {R} for researchers who do not use {R}.
\emph{Advances in Methods and Practices in Psychological Science},
\emph{5}(2). \url{https://doi.org/10.1177/25152459221074654}

\bibitem[\citeproctext]{ref-wickham2023r4ds}
Wickham, H., Çetinkaya-Rundel, M., \& Grolemund, G. (2023). \emph{{R}
for data science} (2nd ed.). O'Reilly Media.
\url{https://r4ds.hadley.nz/}

\end{CSLReferences}


\backmatter


\end{document}
